% =====================================================================
%  Documento independiente: titulo, Resumen y secciones 4 (Materiales y
%  metodos), 5 (Resultados) y 6 (Discusion), con sus referencias, para enviar
%  a revision (Dra. Kruse). Compilar con:  latexmk metodos_resultados_discusion_sol.tex
%  (requiere haber compilado main.tex antes: toma de build/main.aux los numeros
%  de las referencias cruzadas a otras secciones).
%  20 oct. 2026: Sol pidio agregar el titulo y el Resumen (ultima version) y
%  sacar el Anexo A.
% =====================================================================
\documentclass[12pt,a4paper,oneside]{report}
% =====================================================================
%  preambulo.tex — Configuración general del PFC
%  Compilar con: latexmk  (pdflatex + biber; ver .latexmkrc)
% =====================================================================

% ---------- Idioma y codificación ----------
\usepackage[T1]{fontenc}
\usepackage[utf8]{inputenc}
\usepackage[english,spanish,es-tabla,es-noshorthands,es-nodecimaldot]{babel} % español = idioma principal (último)
\usepackage{csquotes}

% ---------- Tipografía (la plantilla ITBA usa Noto Serif) ----------
\usepackage{microtype}
\usepackage{newunicodechar}
% Red de seguridad por si se pega texto con caracteres Unicode sueltos
\newunicodechar{µ}{\textmu}
\newunicodechar{−}{\ensuremath{-}}
\newunicodechar{≤}{\ensuremath{\leq}}
\newunicodechar{≥}{\ensuremath{\geq}}
\newunicodechar{⁺}{\ensuremath{^{+}}}
\newunicodechar{→}{\ensuremath{\rightarrow}}
\newunicodechar{η}{\ensuremath{\eta}}
\newunicodechar{λ}{\ensuremath{\lambda}}

% ---------- Página ----------
\usepackage[a4paper,top=2.5cm,bottom=2.5cm,left=2.5cm,right=2.5cm]{geometry}
\usepackage{setspace}
\onehalfspacing
\setlength{\parskip}{0.5em}
\setlength{\parindent}{0pt}
\setlength{\emergencystretch}{2em}  % evita renglones desbordados

% ---------- Matemática, unidades, tablas, figuras ----------
\usepackage{amsmath}
\usepackage{notomath}   % Noto Serif para texto y matemática (coherente con la plantilla)  % cargar después de amsmath; ya provee los símbolos de amssymb
\usepackage[locale=DE,per-mode=symbol]{siunitx} % DE → coma decimal
\usepackage{graphicx}
\usepackage{tikz}
\usepackage{booktabs}
\usepackage{tabularx}
\usepackage{array}
\usepackage{multirow}
\usepackage{pifont}          % \ding{173} = ②
\usepackage[font=small,labelfont=bf,labelsep=period]{caption}
\usepackage{subcaption}
\renewcommand{\thesubfigure}{\Alph{subfigure}} % paneles (A), (B), ... igual que en los pies de figura
\usepackage{placeins} % \FloatBarrier: evita que las figuras de una seccion floten a la siguiente
\usepackage{float}    % [H]: figura exactamente donde esta en el texto (se usa en el cap. 4)
\usepackage{enumitem}
\usepackage[dvipsnames]{xcolor}
\newcolumntype{C}{>{\centering\arraybackslash}X}
\graphicspath{{figuras/}{logos/}}

% ---------- Títulos: "1. Introducción", "1.1. ..." (sin "Capítulo") ----------
\usepackage{titlesec}
\titleformat{\chapter}[hang]{\normalfont\LARGE\bfseries}{\thechapter.}{0.5em}{}
\titlespacing*{\chapter}{0pt}{0pt}{1.2em}
\titleformat{\section}{\normalfont\Large\bfseries}{\thesection.}{0.5em}{}
\titleformat{\subsection}{\normalfont\large\bfseries}{\thesubsection.}{0.5em}{}
\setcounter{secnumdepth}{3}
\setcounter{tocdepth}{2}

% ---------- Encabezados y pie ----------
\usepackage{fancyhdr}
\pagestyle{fancy}
\fancyhf{}
\fancyfoot[C]{\thepage}
\renewcommand{\headrulewidth}{0pt}
\fancypagestyle{plain}{\fancyhf{}\fancyfoot[C]{\thepage}\renewcommand{\headrulewidth}{0pt}}

% ---------- Bibliografía IEEE ----------
% language=english: las referencias siguen la convención IEEE en inglés
% ("vol.", "no.", "pp.", "Accessed:", "[Online]. Available:"), como en la versión GDocs.
\usepackage[backend=biber,bibstyle=ieee,citestyle=numeric-comp,sorting=none,dashed=false,urldate=long,language=english,autolang=other]{biblatex}
\addbibresource{referencias.bib}
\DeclareSourcemap{\maps[datatype=bibtex]{\map[overwrite=false]{\step[fieldset=langid, fieldvalue=english]}}}
% english.lbx no activa \uspunctuation (coma dentro de las comillas de
% titulo, convencion IEEE); mapear a american.lbx, que si lo hace.
\DeclareLanguageMapping{english}{american}
% @online sin fecha: evitar el "()" vacío
\usepackage{xpatch}
\xpatchbibdriver{online}{\printtext[parens]{\usebibmacro{date}}}{\iffieldundef{year}{}{\printtext[parens]{\usebibmacro{date}}}}{}{}
\DeclareFieldFormat{sentencecase}{#1}   % respetar mayúsculas de los títulos tal como están en el .bib
\renewcommand*{\bibfont}{\small}
\appto{\bibsetup}{\raggedright}    % evita URLs "estiradas" en líneas justificadas
\setcounter{biburlnumpenalty}{9000}   % permite cortar URLs largas
\setcounter{biburllcpenalty}{7000}
\setcounter{biburlucpenalty}{8000}

% ---------- Hipervínculos (siempre al final) ----------
\usepackage[hidelinks,unicode]{hyperref}

% =====================================================================
%  Macros del proyecto
% =====================================================================
% Marcadores de trabajo pendiente — se ven en rojo/naranja en el PDF y se
% buscan fácil con:  grep -rn "pendiente\|confirmar\|revisar" capitulos/
% Para la versión final: \renewcommand{\pendiente}[1]{} etc. (o resolverlos).
\newcommand{\pendiente}[1]{\textcolor{red}{\textbf{[PENDIENTE:} #1\textbf{]}}}
\newcommand{\confirmar}[1]{\textcolor{orange!85!black}{\textbf{[CONFIRMAR:} #1\textbf{]}}}
\newcommand{\revisar}[1]{\textcolor{Purple}{\textbf{[REVISAR:} #1\textbf{]}}}
% Cambios de la última ronda de correcciones: en el PDF no se ven; en el Word
% (herramientas/word/resaltar.py) salen resaltados en amarillo, como pidió la
% Dra. Kruse (23 oct. 2026). Para la versión final se pueden borrar las marcas.
\newcommand{\resaltar}[1]{#1}

% Figura con respaldo: si el archivo no está, dibuja un recuadro con el nombre.
% Uso: \incluirfigura[width=0.8\linewidth]{Fig4_esquema_SLR.png}
\newcommand{\incluirfigura}[2][width=\linewidth]{%
  \IfFileExists{figuras/#2}%
    {\includegraphics[#1]{#2}}%
    {\fbox{\parbox[c][4cm][c]{0.8\linewidth}{\centering\small\ttfamily
      FALTA ARCHIVO: figuras/\detokenize{#2}}}}}

% Espacio extra en el índice para "Anexo A." (se llama desde 9_anexos.tex)
\makeatletter
\newcommand{\anexotocspacing}{\renewcommand*\l@section{\@dottedtocline{1}{1.5em}{5.2em}}}
% Índices de figuras y de tablas: más ancho para los números de dos cifras
% ("5.10.", "5.17."), que quedaban pegados al título (21 oct. 2026)
\renewcommand*\l@figure{\@dottedtocline{1}{1.5em}{2.9em}}
\renewcommand*\l@table{\@dottedtocline{1}{1.5em}{2.9em}}
% Índice general: menos espacio antes de cada capítulo (1,0em -> 0,55em) para que
% la última línea ("Referencias bibliográficas") no quede sola en una página
\patchcmd{\l@chapter}{\vskip 1.0em}{\vskip 0.55em}{}{\PackageWarning{tesis}{no se pudo ajustar el indice general}}
\makeatother

% Abreviaturas de uso frecuente
\newcommand{\stevia}{\textit{Stevia rebaudiana} Bertoni}
\newcommand{\Srebaudiana}{\textit{S.~rebaudiana}}
\newcommand{\anymaze}{AnyMaze\textsuperscript{\textcopyright}}

\usepackage{xr}
\externaldocument{build/main}

\begin{document}
\pagenumbering{arabic}

\begin{center}
  {\Large\textbf{Efectos neurobiológicos a largo plazo de la exposición a \textit{Stevia rebaudiana} durante la etapa juvenil y la adultez}}\\[0.8em]
  {\large Resumen, Materiales y métodos, Resultados, Discusión, Perspectivas futuras y Conclusiones}\\[0.6em]
  Thiago Sayegh --- IByME-CONICET\\[0.3em]
  \textit{Versión para revisión de la Dra. M.\,S. Kruse --- 21 de octubre de 2026}
\end{center}

\vspace{0.8em}
\noindent\textbf{Cómo leer este documento.} Son el Resumen y las secciones 4 a 8 de la tesis, seguidos de su propia lista de referencias (solo las que se citan en estas secciones, numeradas por orden de aparición, por lo que los números difieren de los de la tesis completa). Las referencias a otras secciones conservan la numeración de la tesis. Lo que falta o hay que verificar está marcado en color: \pendiente{falta redactar o conseguir un dato} y \confirmar{dato a confirmar con la tutora}.

% ---------------- Resumen (media carilla, con palabras clave) ----------------
% Plantilla ITBA: "Sintesis breve del problema, los objetivos, la metodologia y los
% resultados principales, acompanada de palabras clave" (media carilla, ~150-180 palabras).
% Version condensada (8 oct. 2026) a partir de la version de la Dra. Kruse (6 oct. 2026,
% disponible en la etiqueta git tesina-20261006); actualizada el 16 oct. 2026 con los
% Resultados corregidos por la Dra. Kruse (consumo de liquidos y neurogenesis). Incluye el
% analisis multivariado de la SLR (16 oct. 2026).
\chapter*{Resumen}
\addcontentsline{toc}{chapter}{Resumen}

El consumo de bebidas endulzadas con edulcorantes ha crecido exponencialmente y se conocen poco sus efectos a largo plazo. Este trabajo comparó los efectos del acceso ilimitado a una bebida endulzada con \stevia{}, con sacarosa o con agua, durante la etapa juvenil (PD25) o la adultez (PD75) de ratas macho Sprague-Dawley. Tras 25 días de exposición y 25 días de acceso exclusivo a agua, se evaluaron la conducta tipo ansiosa (\textit{open field}), la memoria de reconocimiento (NOR) y la separación de patrones espaciales (SLR), y la neurogénesis hipocampal mediante inmunocitoquímica para PCNA y DCX.

Los juveniles consumieron, en relación con su peso, más del doble de edulcorante que los adultos. La stevia aumentó la exploración de las zonas centrales del \textit{open field} sin alterar la locomoción total. En el NOR, la sacarosa juvenil redujo la discriminación del objeto novedoso en todas las demoras, y la stevia lo hizo solo a las 2\,h en adultos y a las 24\,h en juveniles. En la SLR, la sacarosa redujo el índice de discriminación en juveniles (d-SLR y s-SLR) y en adultos (s-SLR), mientras que la stevia no difirió del control. La sacarosa redujo, además, los neuroblastos en proliferación (PCNA$^{+}$/DCX$^{+}$) del giro dentado. Un análisis multivariado de la SLR (PCA, clustering de Ward y Random Forest) mostró que ningún grupo definió un perfil conductual separado y que la stevia se diferenció del control por su exploración, no por su discriminación.

\vspace{1em}
\textbf{Palabras clave:} \textit{Stevia rebaudiana}, sacarosa, neurodesarrollo, separación de patrones, neurogénesis hipocampal, edulcorantes no calóricos.


\clearpage
\setcounter{chapter}{3}
% =====================================================================
\chapter{Materiales y métodos}\label{cap:metodos}
% =====================================================================
% Reescrito el 8 oct. 2026 segun las indicaciones de la Dra. Kruse. Fuentes:
%  - Animales: Coirini et al. 2022 (Biomedicines), seccion 2.1, traducida, mas el grupo stevia.
%  - Extracto de stevia, procedimientos conductuales, OF, SLR, inmunocitoquimica,
%    recuento celular y estadistica: borrador del laboratorio "mat meth draft".
%  - NOR: Coirini et al. 2022, seccion 2.4, traducida.
%  - SLR: texto del borrador (escrito para raton) con las medidas para rata.
% Todo lo estadistico va en la seccion de analisis estadistico, no en los demas apartados.

% ---------------------------------------------------------------------
\section{Animales de experimentación}\label{sec:diseno}

Los procedimientos con animales fueron aprobados por el Comité Institucional de Cuidado y Uso de Animales de Laboratorio (CICUAL) del Instituto de Biología y Medicina Experimental (IByME-CONICET, Argentina; protocolo 19/2025), de acuerdo con los lineamientos para el cuidado y uso de animales de laboratorio del Consejo de la Comunidad Europea (86/609/CEE) y del Instituto Nacional de Salud de los Estados Unidos (NIH). Se utilizaron ratas macho Sprague-Dawley, alojadas en grupos de 3 a 4 animales por caja y mantenidas en un ciclo de 12\,h de luz y 12\,h de oscuridad, con alimento balanceado (Gepsa, Grupo Pilar S.A., Argentina) y agua disponibles \textit{ad libitum}.

Durante un período de 25 días, ratas de 25 días de edad (PD25; grupo juvenil) y ratas de 75 días de edad (PD75; grupo adulto) tuvieron acceso irrestricto a agua (grupos control) o la opción de beber, además de agua, una solución de sacarosa al 10\,\% (p/v; grupos sacarosa) o un extracto acuoso de \stevia{} al 4\,\% (v/v; grupos stevia) \textit{ad libitum}. Las soluciones se renovaron día por medio. Finalizado este período de 25 días de exposición, se retiraron las botellas con la solución de sacarosa o de stevia y todos los animales bebieron únicamente agua durante otros 25 días (Figura~\ref{fig:linea-temporal}) \cite{kruse2019sucrose,coirini2022sucrose}. Al finalizar el protocolo, los animales fueron evaluados en las pruebas conductuales: \textit{open field} (OF), reconocimiento de objeto novedoso (NOR) y reconocimiento espontáneo de posición (SLR). Luego de las pruebas, un subgrupo de animales de los grupos control y sacarosa fue perfundido para los estudios de inmunocitoquímica (\S\ref{sec:histologia}).

El consumo de agua, de solución endulzada (sacarosa o stevia) y de alimento se registró día por medio durante todo el período experimental, y el peso corporal, dos veces por semana. El consumo se midió por caja y se normalizó al peso corporal total de los animales alojados en ella (mL o g por g de peso corporal por día); por lo tanto, la caja se consideró la unidad experimental en todos los análisis de consumo.

\begin{figure}[htbp]
\centering
\begin{tikzpicture}[
    font=\small,
    >=stealth,
    line/.style={draw, black, thick},
    dayline/.style={draw, black!55, thin},
  ]
  % ---- fila juvenil ----
  \node[anchor=east] at (-0.15,2.7) {\textbf{Juvenil}};
  \fill[black!78] (0,2.4) rectangle (3,3.0);
  \fill[black!12] (3,2.4) rectangle (6,3.0);
  \draw[line] (0,2.4) rectangle (6,3.0);
  \draw[line] (3,2.4) -- (3,3.0);
  \foreach \x/\lbl in {0/PD25, 3/PD50, 6/PD75}
    \draw[dayline] (\x,2.4) -- (\x,3.35) node[above, black] {\lbl};

  % ---- fila adulta ----
  \node[anchor=east] at (-0.15,0.3) {\textbf{Adulto}};
  \fill[black!78] (0,0.0) rectangle (3,0.6);
  \fill[black!12] (3,0.0) rectangle (6,0.6);
  \draw[line] (0,0.0) rectangle (6,0.6);
  \draw[line] (3,0.0) -- (3,0.6);
  \foreach \x/\lbl in {0/PD75, 3/PD100, 6/PD125}
    \draw[dayline] (\x,-0.35) -- (\x,0.0) node[below, black] {\lbl};

  % ---- convergencia hacia el desenlace común ----
  \draw[black, thick] (6,2.7) -- (7.1,2.7) -- (7.1,0.3) -- (6,0.3);
  \draw[black, thick, ->] (7.1,1.5) -- (7.9,1.5);
  \node[draw=black, align=center, fill=white, text width=3.5cm, inner sep=6pt] at (10.0,1.5)
    {Pruebas conductuales\\ \footnotesize OF, NOR y SLR\\[2pt] + inmunocitoquímica\\ \footnotesize (PCNA y DCX)};

  % ---- leyenda (apilada) ----
  \fill[black!78] (0,-1.05) rectangle (0.5,-0.7);
  \draw[line] (0,-1.05) rectangle (0.5,-0.7);
  \node[anchor=west] at (0.65,-0.875) {Exposición (25 d): agua, sacarosa o stevia};
  \fill[black!12] (0,-1.65) rectangle (0.5,-1.3);
  \draw[line] (0,-1.65) rectangle (0.5,-1.3);
  \node[anchor=west] at (0.65,-1.475) {Solo agua (25 d)};
\end{tikzpicture}
\caption[Línea temporal de los procedimientos experimentales]{Línea temporal de los procedimientos experimentales. Durante 25 días, ratas PD25 (grupo juvenil) y PD75 (grupo adulto) tuvieron acceso irrestricto a agua (control) o la opción de beber, además de agua, una solución de sacarosa al 10\,\% o un extracto de stevia al 4\,\%. Luego, todos los animales bebieron únicamente agua durante otros 25 días, tras lo cual se realizaron las pruebas conductuales y los estudios de inmunocitoquímica.}
\label{fig:linea-temporal}
\end{figure}

% ---------------------------------------------------------------------
\FloatBarrier
\section{Preparación y caracterización de la infusión de stevia}\label{sec:soluciones}

Las hojas de \stevia{} se obtuvieron de Caña Dulce\textsuperscript{\textregistered} (San Ignacio, Misiones, Argentina), con certificación de Argencert, Argentina (R.P.A.A. N.º~MI000348 y R.P.E.A. N.º~MI000097). El extracto de stevia se preparó según un método descrito previamente \cite{kruse2025stevia}, adaptado de Sharma et~al.\ \cite{sharma2009stevia}, elegido por su alto rendimiento de glucósidos de esteviol \cite{celaya2022infusions} y por su similitud con los procedimientos de preparación orientados a la salud. Se hirvieron 20\,g de hojas de stevia en 1\,L de agua durante 10\,min, se dejó reposar 20\,min, se filtró y se volvió a hervir durante 45--50\,min hasta reducir el volumen a 300\,mL. El extracto se fraccionó en alícuotas de 40\,mL, que se almacenaron a \SI{-20}{\celsius}; las alícuotas se descongelaron día por medio y se diluyeron a 1\,L (4\,\% v/v). La solución de sacarosa al 10\,\% (p/v; Ledesma, Argentina) se preparó en agua. Estas concentraciones se eligieron sobre la base de estudios previos y del consumo voluntario de los animales \cite{coirini2022sucrose,rey2024metformin,kruse2019sucrose,kruse2025stevia}.

La caracterización química del extracto de stevia fue realizada por el Laboratorio Central de la Facultad de Ciencias Exactas, Químicas y Naturales (FCEQyN) de la Universidad Nacional de Misiones (Argentina), siguiendo protocolos establecidos \cite{yildizozturk2015extraction,celaya2022infusions}. El contenido de glucósidos de esteviol se determinó por cromatografía líquida de alta eficacia (HPLC). Los compuestos fenólicos totales y los taninos se cuantificaron con un espectrofotocolorímetro y la actividad antioxidante se evaluó mediante el ensayo de captación del radical 2,2-difenil-1-picrilhidrazilo (DPPH). La composición del extracto se resume en la Tabla~\ref{tab:extracto}.

\begin{table}[htbp]
\centering
\caption[Composición química del extracto de stevia]{Composición química del extracto de stevia. Los valores se expresan como media $\pm$ SEM ($n = 3$). DPPH: 2,2-difenil-1-picrilhidrazilo; GAE: equivalentes de ácido gálico; QE: equivalentes de quercetina; IC$_{50}$: concentración inhibitoria media.}
\label{tab:extracto}
\small
\begin{tabular}{@{}l l@{}}
\toprule
\textbf{Componente} & \textbf{Valor} \\
\midrule
Esteviósido & $3{,}80 \pm 0{,}29$\,mg/mL \\
Rebaudiósido A & $1{,}47 \pm 0{,}10$\,mg/mL \\
Rebaudiósido C & $0{,}25 \pm 0{,}02$\,mg/mL \\
Glucósidos de esteviol totales & $5{,}52 \pm 0{,}42$\,mg/mL \\
Fenoles & $2{,}99 \pm 0{,}38$\,mg GAE/mL de extracto \\
Flavonoides & $1{,}49 \pm 0{,}17$\,mg QE/mL de extracto \\
Taninos & $1{,}73 \pm 0{,}11$\,mg GAE/mL de extracto \\
Actividad antioxidante (DPPH), IC$_{50}$ & $14{,}04 \pm 0{,}56$\,\si{\micro\gram}/mL \\
\bottomrule
\end{tabular}
\end{table}

% ---------------------------------------------------------------------
\FloatBarrier
\section{Procedimientos conductuales}\label{sec:conductual}

Los dispositivos de prueba se iluminaron con una luz de 42\,W suspendida a aproximadamente 1\,m de altura y dirigida hacia el techo. Los animales se aclimataron a la sala de comportamiento durante 20\,min antes de las pruebas, que se realizaron por la mañana (10:00\,h). Los dispositivos se limpiaron con etanol al 70\,\% entre animales. Las sesiones de comportamiento se grabaron en video, se codificaron y se analizaron con el software de seguimiento \anymaze{} (Stoelting Co., EE.\,UU.).

\FloatBarrier
\subsection{\textit{Open field} (OF)}\label{sec:of}

El dispositivo consistió en una arena cuadrada de 75\,cm de lado con paredes opacas de 30\,cm de alto. Las ratas se colocaron en el centro de la arena y se les permitió explorarla libremente durante 5\,min. La arena se dividió digitalmente en tres zonas concéntricas: central, intermedia y periférica (zonas 1 a 3; Figura~\ref{fig:of-esquema}). El piso de la arena está dividido en una cuadrícula de $5 \times 5$ cuadrados de 15\,cm de lado: la zona 1 corresponde al cuadrado central ($15 \times 15$\,cm); la zona 2, al anillo de ocho cuadrados que lo rodea (límite externo de $45 \times 45$\,cm); y la zona 3, a la franja de 15\,cm de ancho adyacente a las paredes. Se registraron la distancia total recorrida, el número de zonas visitadas y las entradas, las distancias recorridas y el tiempo de permanencia en cada zona \cite{kruse2019sucrose}. Se registró además el \textit{rearing}, una conducta de exploración vertical en la que el animal se incorpora sobre las patas traseras \cite{lever2006rearing}, como número de eventos y tiempo total.

\begin{figure}[H]
  \centering
  \incluirfigura[width=0.36\linewidth]{Fig4_OF_zonas.png}
  \caption[Aparato de \textit{open field} (OF) y zonas de registro]{Aparato de \textit{open field} (OF) y zonas de registro: captura cenital de \anymaze{} de la arena (75\,cm de lado), con la cuadrícula del piso (cuadrados de 15\,cm de lado) y las tres zonas definidas digitalmente: central (1), intermedia (2) y periférica (3).}
  \label{fig:of-esquema}
\end{figure}

\FloatBarrier
\subsection{Reconocimiento de objeto novedoso (NOR)}\label{sec:nor}

El aparato consistió en una caja abierta de madera oscura (75\,cm de largo $\times$ 30\,cm de alto $\times$ 55\,cm de ancho). Los objetos a discriminar diferían en forma y textura; eran de vidrio, plástico o cerámica y tenían cuatro formas distintas (frascos de vidrio, una rana de juguete de plástico, cajas de plástico y tazas de cerámica), y no podían ser desplazados por las ratas \cite{perezgarcia2016malnutrition,leger2013nor}. Brevemente, durante la semana previa a la prueba las ratas fueron manipuladas una vez al día durante 3 días consecutivos. Tres días antes de la prueba se permitió que las ratas exploraran el aparato durante 5\,min por día. En estas sesiones de habituación no había objetos en la caja. Veinticuatro horas después de la última sesión de habituación, las ratas fueron entrenadas para el reconocimiento de objetos permitiéndoles explorar durante 5\,min dos objetos idénticos colocados en la arena de prueba (fase de muestreo). Durante esta fase de muestreo (T1), los objetos (dos frascos de vidrio de 12,5\,cm de alto y 5\,cm de diámetro) se colocaron al azar en dos esquinas opuestas del aparato, a 10\,cm de las paredes laterales (Figura~\ref{fig:nor-objetos}A). La rata se colocó en el centro del aparato y se le permitió explorar los dos objetos idénticos. Después de T1, la rata fue devuelta a su caja y, tras un intervalo entre ensayos de 2\,h (memoria a corto plazo, T2), de 24\,h (memoria a largo plazo, T3) y de 7 días desde la fase de muestreo (memoria consolidada, T4), fue evaluada nuevamente en los ensayos de elección. Durante T2, T3 y T4, un objeto novedoso reemplazó a uno de los objetos presentados durante T1. Así, las ratas fueron reexpuestas durante 5\,min a dos objetos: una copia del objeto familiar (F) y el objeto novedoso (N) (Figura~\ref{fig:nor-objetos}B); es decir, un frasco de vidrio se reemplazó por una rana de juguete de plástico, por una caja de plástico anaranjada de 10,5\,cm de largo $\times$ 9\,cm de ancho $\times$ 4\,cm de alto y por una taza cónica de cerámica blanca de 8,5\,cm de alto y 8,5\,cm de diámetro en T2, T3 y T4, respectivamente (Figura~\ref{fig:nor-esquema}; los objetos se muestran en la Figura~\ref{fig:nor-tasa}). Todas las combinaciones y ubicaciones de los objetos se contrabalancearon para reducir el posible sesgo causado por la preferencia por ubicaciones u objetos particulares. Las sesiones T1 a T4 se grabaron en video y se analizaron los cruces de línea con el software \anymaze{} (Stoelting Co., Wood Dale, IL, EE.\,UU.). El tiempo que las ratas dedicaron a explorar cada objeto durante T1 a T4 fue registrado manualmente por un único experimentador que desconocía la asignación de los grupos. La exploración se definió de la siguiente manera: dirigir la nariz hacia el objeto a una distancia de 2\,cm o menos y/o tocar el objeto con la nariz. Dar vueltas alrededor del objeto o sentarse sobre él no se consideró comportamiento exploratorio. Los datos se informan como la tasa de exploración, que es el tiempo dedicado a explorar el objeto novedoso ($t_N$) dividido por el tiempo dedicado a explorar ambos objetos:
\begin{equation}
  \text{tasa de exploración} = \frac{t_N}{t_N + t_F}.
  \label{eq:tasa}
\end{equation}

\begin{figure}[htbp]
  \centering
  \incluirfigura[width=0.9\linewidth]{Fig4_NOR_esquema.png}
  \caption[Esquema del protocolo de la prueba NOR]{Esquema del protocolo de la prueba NOR. En la fase de muestreo (T1) se presentan dos objetos idénticos (A); en los ensayos de elección, realizados 2\,h (T2), 24\,h (T3) y 7 días (T4) después de T1, uno de ellos se reemplaza por un objeto novedoso distinto en cada ensayo (B, C y D). Esquema cedido por la Dra. M.\,S. Kruse.}
  \label{fig:nor-esquema}
\end{figure}

\begin{figure}[htbp]
  \centering
  \begin{subfigure}[b]{0.38\linewidth}
    \centering
    \incluirfigura[width=\linewidth]{Fig4_NOR_objetos_iguales_recortada.png}
    \caption{}
  \end{subfigure}
  \hspace{1.5em}
  \begin{subfigure}[b]{0.38\linewidth}
    \centering
    \incluirfigura[width=\linewidth]{Fig4_NOR_objetos_diferentes_recortada.png}
    \caption{}
  \end{subfigure}
  \caption[Fases del NOR]{Fases del NOR en vista cenital. (A)~Fase de muestreo (T1): dos objetos idénticos. (B)~Ensayos de elección (T2 a T4): una copia del objeto familiar y un objeto novedoso.}
  \label{fig:nor-objetos}
\end{figure}

\FloatBarrier
\subsection{Reconocimiento espontáneo de posición (SLR)}\label{sec:slr}

Este procedimiento se realizó como fue descrito por Reichelt et~al.\ \cite{reichelt2016sucrose,reichelt2021slr}. El aparato consistió en un campo abierto circular de plástico gris y base opaca, de 100\,cm de diámetro $\times$ 45\,cm de alto, sin techo. La arena estuvo rodeada de tres señales espaciales distintivas (pósteres negros y blancos con patrones únicos) de 35\,cm de alto $\times$ 50\,cm de ancho, ubicadas equidistantes alrededor de la arena, y de las señales distales proporcionadas por los muebles de la sala, incluidos los estantes y la puerta.

Durante la semana previa a la habituación, los animales fueron manipulados una vez al día durante 5 días consecutivos. Antes del día del primer ensayo se permitió que cada animal explorara el dispositivo durante 5\,min por día durante 3 días consecutivos (habituación). Los objetos fueron tazas con dibujos o ranas de juguete (Figura~\ref{fig:objetos-slr}), limpiados con etanol al 70\,\% entre pruebas; dentro de cada sesión, los objetos de la fase de muestreo y de la fase de prueba fueron copias idénticas.

En el primer ensayo se expuso a las ratas a tres objetos idénticos, A1, A2 y A3, durante 10\,min (fase de muestreo). La disposición de los objetos dependió de la configuración evaluada: d-SLR o s-SLR (Figura~\ref{fig:esquema-slr}). En ambas configuraciones, los objetos se colocaron a lo largo de una circunferencia, a 10\,cm de la pared. En cada experimento, la mitad de los animales fue asignada a la tarea d-SLR y la otra mitad a la s-SLR ($n$ = 9--10 por grupo y configuración; \ref{anx:muestra}). Para la d-SLR, los objetos A1, A2 y A3 se colocaron equidistantes, a 120° y a una distancia aproximada de 49\,cm entre ellos. Para la s-SLR, los objetos A2 y A3 se colocaron a 50° de distancia, a 20,5\,cm entre ellos, mientras que A1 se ubicó a la misma distancia de A2 y de A3. El ensayo se fundamenta en la separación de patrones: la demanda cognitiva es mayor cuanto menor es la distancia entre los objetos, de modo que la s-SLR impone una demanda cognitiva mayor, y es más dependiente de la neurogénesis, que la d-SLR \cite{reichelt2021slr}. Veinticuatro horas después de la fase de muestreo, los animales regresaron a la arena, donde fueron expuestos solo a dos objetos, uno familiar (A1) y otro novedoso (N), durante 5\,min (fase de prueba). El objeto N se ubicó en la posición media entre las que ocupaban los objetos A2 y A3 en la fase de muestreo (Figura~\ref{fig:esquema-slr}).

\begin{figure}[htbp]
  \centering
  \incluirfigura[width=0.9\linewidth]{Fig4_esquema_SLR_demanda.png}
  \caption[Esquema de la prueba de reconocimiento espontáneo de posición (SLR)]{Esquema de la prueba de reconocimiento espontáneo de posición (SLR). En la fase de muestreo (día 1, 10\,min) el animal explora tres objetos idénticos (A1, A2 y A3) en la configuración d-SLR (fila superior) o s-SLR (fila inferior). En la fase de prueba (día 2, 24\,h después, 5\,min) se presentan dos objetos: uno en la posición familiar de A1 y otro en una posición novedosa (N), ubicada entre las posiciones que antes ocupaban A2 y A3 (círculos punteados). Los símbolos oscuros alrededor de la arena representan las señales espaciales, y la flecha indica la mayor demanda cognitiva de la configuración s-SLR. Adaptado de \cite{reichelt2021slr}.}
  \label{fig:esquema-slr}
\end{figure}

\begin{figure}[htbp]
  \centering
  \incluirfigura[width=0.7\linewidth]{Fig4_objetos_SLR.jpg}
  \caption{Objetos utilizados en la prueba SLR: tazas con dibujos y ranas de juguete.}
  \label{fig:objetos-slr}
\end{figure}

Los datos se informan como el índice de discriminación (DI), que es la diferencia entre el tiempo dedicado a explorar el objeto en la posición novedosa ($t_N$) y el tiempo de exploración del objeto en la posición familiar ($t_F$), dividida por el tiempo total de exploración de ambos objetos:
\begin{equation}
  \mathrm{DI} = \frac{t_N - t_F}{t_N + t_F}.
  \label{eq:di}
\end{equation}
Se informa también la tasa de exploración equivalente (Ecuación~\ref{eq:tasa}), de la cual el DI es una transformación lineal ($\mathrm{DI} = 2\cdot\text{tasa} - 1$). En la fase de muestreo se registraron, además, el tiempo de exploración de cada uno de los tres objetos y la distancia recorrida.

\FloatBarrier
\subsection{Registro con \anymaze}\label{sec:principio}

El registro se realizó con el software de videoseguimiento \anymaze{} (Stoelting Co., EE.\,UU.; versión 7.66), que opera sobre el video cenital de cada sesión. El programa segmenta al animal del fondo por contraste y sigue su posición cuadro a cuadro, reconstruyendo la trayectoria del centroide. De esa trayectoria deriva de manera automática un conjunto de variables cinemáticas y de ocupación: la distancia total recorrida, la velocidad media, los cruces de línea (\textit{line crossings}) sobre una grilla virtual y, mediante zonas definidas por el usuario ---regiones de interés trazadas alrededor de cada objeto o sector de la arena---, el número de entradas y el tiempo de permanencia en cada zona.

Permanecer dentro de la zona de un objeto no equivale a explorarlo: el seguimiento detecta posición, no atención. Un animal que se sienta sobre un objeto, se acicala a su lado o descansa dentro de la zona acumularía \enquote{tiempo en zona} sin exploración real. Por eso la exploración activa ---definida como dirigir la nariz hacia el objeto a 2\,cm o menos o tocarlo con la nariz, excluyendo treparlo o posarse sobre él--- se cuantificó de forma manual mediante la función de marcación por tecla de \anymaze: revisando el video, un único experimentador ciego al tratamiento mantuvo presionada una tecla mientras el animal ejecutaba conducta exploratoria hacia un objeto dado, y el software integró ese tiempo por objeto. El \textit{rearing} del OF se puntuó de la misma manera. Por esta razón, el registro de cada prueba contiene dos familias de variables: la exploración activa, de origen manual, de la que se calculan la tasa de exploración y el índice de discriminación; y las variables cinemáticas y de ocupación, de origen automático, que se emplean como control locomotor y como lectura de conducta tipo ansiosa en el OF. Todas las sesiones se grabaron, se codificaron y se cuantificaron de forma diferida, no en vivo, de modo que el experimentador desconocía el tratamiento de cada animal y el criterio se mantuvo uniforme.

% ---------------------------------------------------------------------
\FloatBarrier
\section{Inmunocitoquímica y microscopía de fluorescencia (IHC)}\label{sec:histologia}

Los animales fueron anestesiados profundamente mediante una inyección intraperitoneal de ketamina-xilacina (120 + 12\,mg/kg; Holliday/Richmond, Argentina) y fijados por perfusión intracardíaca con 600\,mL de paraformaldehído (PFA) al 4\,\% frío en PBS (pH~7,4) \cite{kruse2012lxr}. Los cerebros se extrajeron de inmediato y se dejaron en PFA al 4\,\% durante toda la noche a \SI{4}{\celsius}. Luego se lavaron con PBS y se seccionaron con un vibrátomo. Se obtuvieron cortes coronales de \SI{70}{\micro\metre} de espesor (Bregma $-2{,}8$ a $-4{,}2$\,mm), determinados con un atlas del cerebro de rata \cite{paxinos1986rat}, que se incubaron en PBS con 0,1\,\% de Tween~20 y 7\,\% de suero normal de burro durante 1\,h a temperatura ambiente. Las muestras de tejido se incubaron luego con un anticuerpo de ratón anti-PCNA (dilución 1:250; ref.\ MAB424, Merck) y un anticuerpo de cabra anti-DCX (dilución 1:200; ref.\ sc-8066, Santa Cruz Biotechnology) en PBS con 2\,\% de suero de burro y 0,1\,\% de Tween~20, durante toda la noche a \SI{4}{\celsius}. Posteriormente se lavaron en PBS durante 30\,min y se incubaron con un anticuerpo de burro anti-ratón Alexa~488 y un anticuerpo de burro anti-cabra Alexa~555 (ref.\ A21202 y A21432, respectivamente; Thermo Fisher; ambos en dilución 1:1000) durante 2\,h a temperatura ambiente. Los núcleos se contratiñeron con DAPI (1\,\si{\micro\gram}/mL). Finalmente, después de los lavados, los cortes se montaron en portaobjetos de vidrio y se examinaron con un microscopio confocal de disco giratorio Olympus IX81-DSU. Todas las imágenes se obtuvieron en las mismas condiciones: el tiempo de exposición, la intensidad de la lámpara y los filtros se mantuvieron constantes. En algunos cortes se omitió el anticuerpo primario como control; estos cortes se procesaron con el mismo protocolo descrito, y la intensidad de fluorescencia obtenida en ellos se utilizó como referencia para identificar la marca positiva. Todos los portaobjetos se codificaron y fueron analizados por dos observadores independientes que desconocían el tratamiento.

\FloatBarrier
\subsection{Recuento celular}\label{sec:recuento}

Los cortes se tomaron en una serie 1:4 (separados por \SI{280}{\micro\metre}; 5 cortes por animal) y se analizaron con ImageJ Fiji win64 (\url{http://imagej.nih.gov/ij}). Las células positivas para PCNA y/o DCX se contaron dentro de un área delimitada que comprendía la zona subgranular del giro dentado. Para el recuento se consideraron tanto la hoja superior como la inferior de la capa de células granulares y de la zona subgranular. El número de células inmunopositivas fue contado manualmente por dos observadores independientes que desconocían el protocolo experimental. Los resultados se expresan como número de células positivas por mm$^2$; en el análisis estadístico, cada corte se consideró como un valor.

% ---------------------------------------------------------------------
\FloatBarrier
\section{Análisis estadístico}\label{sec:estadistico}

El tamaño muestral se estimó sobre la base de estudios previos (95\,\% de confianza, 10\,\% de error y 80\,\% de potencia estadística) \cite{coirini2022sucrose,rey2024metformin,kruse2019sucrose,kruse2025stevia}, siguiendo las recomendaciones estándar \cite{charan2013sample}.

La normalidad de los datos y la homogeneidad de las varianzas se evaluaron mediante las pruebas de Shapiro-Wilk y de Brown-Forsythe, respectivamente. En este estudio se verificó la homogeneidad de las varianzas (prueba de Brown-Forsythe, $p > 0{,}05$), salvo en el recuento de células PCNA$^{+}$/DCX$^{+}$ (\S\ref{sec:res-neurogenesis}), y no se aplicaron correcciones adicionales.

El consumo de líquidos a lo largo del tiempo se analizó mediante ANOVA de medidas repetidas (ANOVA MR) de dos vías, considerando el tratamiento (control vs.\ sacarosa vs.\ stevia) como factor entre sujetos y el tiempo como factor de medidas repetidas. Se empleó ANOVA de dos vías para analizar los parámetros informados en el OF, el NOR, la SLR y la IHC. La exploración de los objetos en las fases de muestreo del NOR (T1) y de la SLR (día 1) se analizó mediante un ANOVA mixto de tres factores, con la edad y el tratamiento como factores entre sujetos y la posición del objeto como factor intrasujeto.

Cuando se detectaron efectos principales estadísticamente significativos, se aplicó la prueba \textit{post hoc} de Tukey para ajustar las comparaciones múltiples (en los análisis de consumo, la prueba de Šídák). Cuando los residuos se apartaron de la normalidad, las diferencias se confirmaron con la prueba de Mann-Whitney. El nivel de significación estadística se estableció en $p < 0{,}05$, y los datos se expresaron como media $\pm$ error estándar de la media (SEM).

Los resultados provienen de dos experimentos independientes, correspondientes a dos cohortes experimentales realizadas en distintos momentos, utilizando grupos independientes de animales, pero siguiendo el mismo diseño experimental y los mismos protocolos. Los datos de ambas cohortes se combinaron para realizar los análisis estadísticos. Todos los análisis se llevaron a cabo con el programa GraphPad Prism, versión 11.1.0, salvo el de la exploración total de la SLR, que se realizó sobre el logaritmo del tiempo en Python (statsmodels y SciPy).

% =====================================================================
\chapter{Resultados}\label{cap:resultados}
% =====================================================================
% Estilo de los articulos del laboratorio y correcciones de la Dra. Kruse (16 oct. 2026):
% en el post hoc solo "control vs. tratado: p = ..." (sin diferencias de medias ni IC en el texto);
% eta^2_p solo en las tablas; sin pruebas contra el nivel de azar; tablas al final de cada seccion.

\section{Consumo de líquidos}\label{sec:res-consumo}

\subsection{Consumo de líquidos durante el tratamiento}\label{sec:res-consumo-tratamiento}

El curso temporal del consumo difirió según la edad y el endulzante (ANOVA MR de dos vías; grupo juvenil: $F_{\text{Tratamiento} \times \text{Tiempo}}(44, 288) = 5{,}341$; $p < 0{,}0001$; grupo adulto: $F_{\text{Tratamiento} \times \text{Tiempo}}(44, 330) = 2{,}339$; $p < 0{,}0001$; Figura~\ref{fig:consumo-liquidos}A y B, respectivamente). Los animales expuestos a sacarosa consumieron más sacarosa que el control (agua) durante todo el período de exposición, en ambas edades. En cambio, los animales expuestos a stevia consumieron más volumen de stevia que el control por un breve lapso en PD27 (grupo juvenil) y en PD89 (grupo adulto).

Promediado a lo largo del período de exposición, los animales juveniles consumieron, en relación con su peso corporal, más solución endulzada que los adultos, tanto de sacarosa (2,4 veces más) como de stevia (2,8 veces más; ANOVA de dos vías, $F_{\text{Edad} \times \text{Tratamiento}}(1, 64) = 63{,}19$; $p < 0{,}0001$; Figura~\ref{fig:consumo-liquidos-promedio}). Es decir, la dosis de sacarosa o de stevia que recibieron los animales juveniles fue más del doble que la de los adultos.

Los animales de ambas edades prefirieron la solución endulzada al agua, y la preferencia fue mayor por la sacarosa que por la stevia, sin diferencias entre edades (índice de preferencia: consumo de solución endulzada dividido por el consumo total de líquido; sacarosa: 0,96 en juveniles y 0,95 en adultos; stevia: 0,75 en juveniles y 0,79 en adultos; ANOVA de dos vías, $F_{\text{Interacción}}(1, 24) = 0{,}56$; $p = 0{,}46$; $F_{\text{Tratamiento}}(1, 24) = 37{,}80$; $p < 0{,}0001$; $F_{\text{Edad}}(1, 24) = 0{,}33$; $p = 0{,}57$; prueba \textit{post hoc} de Tukey, sacarosa vs.\ stevia: juvenil, $p = 0{,}0010$; adulto, $p = 0{,}0043$; Figura~\ref{fig:preferencia}). Dado que los residuos se apartaron de la normalidad (Shapiro-Wilk, $p = 0{,}024$), las diferencias se confirmaron con la prueba de Mann-Whitney (juvenil, $p = 0{,}002$; adulto, $p < 0{,}001$). Por lo tanto, la mayor dosis de los juveniles no se explica por una mayor preferencia por el sabor dulce, sino por un mayor consumo asociado con la edad.

\begin{figure}[htbp]
  \centering
  \incluirfigura[width=0.45\linewidth]{Fig5_indice_preferencia.pdf}
  \caption[Índice de preferencia por la solución endulzada]{Índice de preferencia por la solución endulzada (consumo de solución dividido por el consumo total de líquido, promediado durante el tratamiento) en los grupos expuestos a sacarosa y a stevia, por edad. Las barras indican la media $\pm$ SEM. Rojo: grupo sacarosa; verde: grupo stevia. Los corchetes indican la comparación de Tukey entre tratamientos dentro de cada edad ($^{**}$~$p < 0{,}01$). $n = 7$ cajas por grupo.}
  \label{fig:preferencia}
\end{figure}

\begin{figure}[htbp]
  \centering
  % figura compuesta (herramientas/figuras/componer.py)
  \incluirfigura[width=\linewidth]{compuestas/Fig5_consumo_liquidos.pdf}
  \caption[Curso temporal del consumo de líquidos durante el tratamiento]{Curso temporal del consumo de líquidos durante el tratamiento. (A)~Grupo juvenil. (B)~Grupo adulto. Eje horizontal: PD; eje vertical: consumo de líquido normalizado al peso corporal (mL/g de peso corporal/día). Cada curva corresponde a una condición y a un fluido (agua o solución endulzada); la leyenda es común a ambos gráficos. Datos expresados como media $\pm$ SEM. $n = 7$ cajas por grupo.}
  \label{fig:consumo-liquidos}
\end{figure}

\begin{figure}[htbp]
  \centering
  \incluirfigura[width=0.9\linewidth]{Fig5_consumo_liquidos_promedio_es_fino.pdf}
  \caption[Consumo de líquidos promediado durante el tratamiento]{Consumo de líquidos promediado a lo largo del período de tratamiento, por condición y edad. Datos expresados como media $\pm$ SEM. Los asteriscos indican diferencias entre la exposición juvenil y la adulta dentro de cada condición (prueba de Tukey; $^{*}\,p < 0{,}05$; $^{***}\,p < 0{,}001$; $^{****}\,p < 0{,}0001$). $n = 7$ cajas por grupo.}
  \label{fig:consumo-liquidos-promedio}
\end{figure}

\subsection{Consumo de agua luego del tratamiento con los endulzantes (período agua)}\label{sec:res-consumo-lavado}

Tras la finalización del tratamiento, se monitoreó el consumo de líquidos durante 25 días adicionales (Figura~\ref{fig:agua-lavado}). El ANOVA MR de dos vías reveló interacciones significativas entre edad y tratamiento ($F_{\text{Edad} \times \text{Tratamiento}}(2, 36) = 3{,}61$; $p = 0{,}037$), edad y tiempo ($F_{\text{Edad} \times \text{Tiempo}}(4{,}73;\ 170{,}30) = 8{,}55$; $p < 0{,}0001$) y tratamiento y tiempo ($F_{\text{Tratamiento} \times \text{Tiempo}}(9{,}46;\ 170{,}30) = 6{,}37$; $p < 0{,}0001$), así como una interacción triple significativa entre edad, tratamiento y tiempo ($F_{\text{Edad} \times \text{Tratamiento} \times \text{Tiempo}}(9{,}46;\ 170{,}30) = 7{,}52$; $p < 0{,}0001$).

Al observar los primeros días del período agua, se encontró que los animales expuestos a las bebidas endulzadas durante la etapa juvenil mostraron signos de frustración, ya que consumieron un menor volumen de agua en comparación con los animales del grupo control. Al examinar solamente los primeros nueve días posteriores al tratamiento (D1--D9), el ANOVA MR de dos vías mostró una interacción significativa entre edad, tratamiento y tiempo ($F_{\text{Edad} \times \text{Tratamiento} \times \text{Tiempo}}(4{,}28;\ 76{,}95) = 2{,}64$; $p = 0{,}037$), así como efectos significativos del tratamiento ($F_{\text{Tratamiento}}(2, 36) = 5{,}21$; $p = 0{,}0103$) y de la edad ($F_{\text{Edad}}(1, 36) = 7{,}95$; $p = 0{,}0078$). Durante este período, los animales juveniles previamente expuestos a stevia presentaron una reducción significativa del consumo promedio de líquidos del 24,97\% respecto de su control ($p = 0{,}003$).

Al extender el análisis a los primeros quince días (D1--D15), se observaron interacciones significativas entre edad y tiempo ($F_{\text{Edad} \times \text{Tiempo}}(2{,}92;\ 105{,}26) = 6{,}68$; $p = 0{,}0004$), tratamiento y tiempo ($F_{\text{Tratamiento} \times \text{Tiempo}}(5{,}85;\ 105{,}26) = 3{,}77$; $p = 0{,}0021$) y edad, tratamiento y tiempo ($F_{\text{Edad} \times \text{Tratamiento} \times \text{Tiempo}}(5{,}85;\ 105{,}26) = 7{,}89$; $p < 0{,}0001$). Durante este intervalo, los juveniles previamente expuestos a sacarosa mantuvieron un consumo promedio inferior al del control ($-14{,}59$\%), aunque esta diferencia no alcanzó significancia estadística ($p = 0{,}051$).

Por el contrario, el análisis del período D11--D25 reveló interacciones significativas entre edad y tratamiento ($F_{\text{Edad} \times \text{Tratamiento}}(2, 36) = 10{,}47$; $p = 0{,}00026$), entre tratamiento y tiempo ($F_{\text{Tratamiento} \times \text{Tiempo}}(5{,}59;\ 100{,}55) = 3{,}91$; $p = 0{,}0019$) y entre edad y tiempo ($F_{\text{Edad} \times \text{Tiempo}}(2{,}79;\ 100{,}55) = 2{,}93$; $p = 0{,}041$), así como un efecto significativo del tiempo ($F_{\text{Tiempo}}(2{,}79;\ 100{,}55) = 11{,}05$; $p < 0{,}0001$). En este período, los animales que consumieron stevia durante su juventud mostraron un incremento del consumo, con valores promedio un 24,83\% superiores a los del control ($p = 0{,}00024$). Por su parte, los que consumieron sacarosa durante su juventud mostraron una aproximación progresiva a los valores del control, sin diferencias significativas durante D17--D25 ($p > 0{,}05$). A diferencia de los grupos juveniles, los grupos adultos no presentaron diferencias significativas respecto de su control de edad en todo el período estudiado (D1--D25).

\begin{figure}[htbp]
  \centering
  \incluirfigura[width=0.8\linewidth]{Fig5_agua_periodo_agua_D1_D25_es_control.pdf}
  \caption[Consumo de agua luego del tratamiento con los endulzantes]{Consumo de agua luego del tratamiento con los endulzantes (período agua), expresado como porcentaje del grupo control de la misma edad (línea azul: 100\%), desde el día 1 (D1) hasta el día 25 (D25) posteriores al tratamiento. Datos expresados como media $\pm$ SEM, según la edad y el tratamiento previo; $n = 7$ cajas por grupo.}
  \label{fig:agua-lavado}
\end{figure}

% ---------------------------------------------------------------------
\FloatBarrier
\section{Conducta tipo ansiosa (OF)}\label{sec:res-of}

\label{sec:res-of-locomocion}La actividad locomotora y la exploración de la arena se evaluaron en el OF al finalizar el protocolo. No se encontraron diferencias en la distancia total recorrida entre los seis grupos (ANOVA de dos vías, $F_{\text{Interacción}}(2, 66) = 0{,}06$; $p = 0{,}94$; $F_{\text{Tratamiento}}(2, 66) = 0{,}90$; $p = 0{,}41$; $F_{\text{Edad}}(1, 66) = 0{,}31$; $p = 0{,}58$; Figura~\ref{fig:of-centro}A), indicando que ni la sacarosa ni la stevia modificaron la actividad locomotora general. Sin embargo, los animales expuestos a stevia fueron los que visitaron más zonas, independientemente de la edad de exposición (ANOVA de dos vías, $F_{\text{Interacción}}(2, 63) = 3{,}28$; $p = 0{,}044$; $F_{\text{Tratamiento}}(2, 63) = 37{,}01$; $p < 0{,}0001$; $F_{\text{Edad}}(1, 63) = 0{,}03$; $p = 0{,}87$; prueba \textit{post hoc} de Tukey, control vs.\ stevia: $p < 0{,}0001$ en ambas edades; Figura~\ref{fig:of-centro}B). Contrariamente, los juveniles expuestos a sacarosa fueron los que menos zonas visitaron (control vs.\ sacarosa: $p < 0{,}05$), sin diferencias en el grupo adulto ($p > 0{,}05$).

\label{sec:res-of-centro}La exploración de las zonas centrales también dependió del tratamiento y, en la distancia recorrida, de su interacción con la edad. Los animales expuestos a stevia ingresaron más veces a las zonas centrales (zonas 1+2) y recorrieron en ellas mayor distancia que los otros grupos, independientemente de la edad de exposición (entradas a la zona 1+2: ANOVA de dos vías, $F_{\text{Interacción}}(2, 57) = 1{,}88$; $p = 0{,}16$; $F_{\text{Tratamiento}}(2, 57) = 91{,}79$; $p < 0{,}0001$; $F_{\text{Edad}}(1, 57) = 3{,}51$; $p = 0{,}066$; distancia en la zona 1+2: $F_{\text{Interacción}}(2, 69) = 5{,}05$; $p = 0{,}009$; $F_{\text{Tratamiento}}(2, 69) = 38{,}88$; $p < 0{,}0001$; $F_{\text{Edad}}(1, 69) = 3{,}15$; $p = 0{,}080$; Figura~\ref{fig:of-centro}C y D). En la prueba \textit{post hoc} de Tukey, las entradas a la zona 1+2 fueron mayores con stevia que en el control en ambas edades (control vs.\ stevia: $p < 0{,}0001$), al igual que la distancia allí recorrida (control vs.\ stevia: juvenil, $p < 0{,}01$; adulto, $p < 0{,}0001$). Por el contrario, solo los animales que consumieron sacarosa durante la etapa juvenil ingresaron menos veces a la zona 1+2 y recorrieron en ella menor distancia que sus controles (control vs.\ sacarosa: entradas, $p < 0{,}01$; distancia, $p < 0{,}001$), sin diferencias en el grupo adulto expuesto a sacarosa respecto de su control de edad ($p > 0{,}05$).

De manera similar, los animales expuestos a stevia ingresaron más veces a la zona 1, aunque la diferencia fue significativa solo en el grupo adulto (entradas a la zona 1: ANOVA de dos vías, $F_{\text{Interacción}}(2, 72) = 2{,}29$; $p = 0{,}11$; $F_{\text{Tratamiento}}(2, 72) = 4{,}11$; $p = 0{,}020$; $F_{\text{Edad}}(1, 72) = 6{,}20$; $p = 0{,}015$; control vs.\ stevia en adultos: $p < 0{,}01$; Figura~\ref{fig:of-centro}E). La distancia recorrida en la zona 1, aunque fue mayor para el grupo stevia, no llegó a ser significativa (distancia en la zona 1: $F_{\text{Interacción}}(2, 74) = 0{,}94$; $p = 0{,}39$; $F_{\text{Tratamiento}}(2, 74) = 2{,}80$; $p = 0{,}067$; $F_{\text{Edad}}(1, 74) = 3{,}85$; $p = 0{,}053$; Figura~\ref{fig:of-centro}F).

En conjunto, estos resultados indican que los animales expuestos a stevia presentaron una mayor exploración de las zonas centrales de la arena que sus respectivos controles y que los animales expuestos a sacarosa, sin modificaciones en la distancia total recorrida, e independientemente de la edad del consumo. En contraste, únicamente los animales expuestos a sacarosa durante la etapa juvenil mostraron una menor exploración de las zonas centrales respecto de sus controles, evidenciada por una reducción tanto en el número de entradas como en la distancia recorrida en estas zonas.

\begin{figure}[htbp]
  \centering
  % figura compuesta (herramientas/figuras/componer.py)
  \incluirfigura[width=\linewidth]{compuestas/Fig5_OF_zonas.pdf}
  \caption[Actividad locomotora y exploración de la arena en el OF]{Actividad locomotora y exploración de la arena en el OF. (A)~Distancia total recorrida. (B)~Número de zonas visitadas. (C)~Entradas a la zona 1+2. (D)~Distancia recorrida en la zona 1+2. (E)~Entradas a la zona 1. (F)~Distancia recorrida en la zona 1. Las barras indican la media $\pm$ SEM. Azul: grupo control; rojo: grupo sacarosa; verde: grupo stevia. Los corchetes indican las comparaciones de Tukey significativas entre tratamientos dentro de cada edad ($^{*}$~$p < 0{,}05$; $^{**}$~$p < 0{,}01$; $^{***}$~$p < 0{,}001$; $^{****}$~$p < 0{,}0001$). $n = 7$--$17$ animales por grupo.}
  \label{fig:of-centro}
\end{figure}

\label{sec:res-of-rearing}En la conducta de \textit{rearing}, en la que el animal se levanta sobre sus patas traseras, el número de eventos dependió de la interacción entre el tratamiento y la edad, mientras que el tiempo dedicado a esta tarea dependió del tratamiento y de la edad, sin interacción entre ambos factores (número de eventos: ANOVA de dos vías, $F_{\text{Interacción}}(2, 73) = 6{,}67$; $p = 0{,}0022$; $F_{\text{Tratamiento}}(2, 73) = 12{,}69$; $p < 0{,}0001$; $F_{\text{Edad}}(1, 73) = 3{,}54$; $p = 0{,}064$; tiempo: $F_{\text{Interacción}}(2, 72) = 2{,}24$; $p = 0{,}11$; $F_{\text{Tratamiento}}(2, 72) = 8{,}81$; $p = 0{,}0004$; $F_{\text{Edad}}(1, 72) = 10{,}12$; $p = 0{,}0022$; Figura~\ref{fig:of-rearing}). Entre los juveniles, tanto los animales expuestos a sacarosa como los expuestos a stevia realizaron menos eventos de \textit{rearing} (prueba \textit{post hoc} de Tukey, control vs.\ sacarosa: $p < 0{,}0001$; control vs.\ stevia: $p < 0{,}0001$; Figura~\ref{fig:of-rearing}A) y le dedicaron menos tiempo a la tarea (control vs.\ sacarosa: $p < 0{,}05$; control vs.\ stevia: $p < 0{,}01$; Figura~\ref{fig:of-rearing}B). En el grupo adulto, solo los animales expuestos a stevia presentaron un número menor de eventos que los expuestos a sacarosa (sacarosa vs.\ stevia: $p < 0{,}05$), aunque no difirieron de los controles ($p > 0{,}05$), y dedicaron menos tiempo a la tarea que ambos grupos (control vs.\ stevia y sacarosa vs.\ stevia: $p < 0{,}05$). Estos resultados indican que la exposición juvenil a cualquiera de los dos endulzantes redujo la exploración vertical, mientras que en la adultez solo la stevia tuvo ese efecto.

\begin{figure}[htbp]
  \centering
  % figura compuesta (herramientas/figuras/componer.py)
  \incluirfigura[width=\linewidth]{compuestas/Fig5_OF_rearing_compuesta.pdf}
  \caption[Conducta de \textit{rearing} en el OF]{Conducta de \textit{rearing} en el OF. (A)~Número de eventos. (B)~Tiempo. Las barras indican la media $\pm$ SEM. Azul: grupo control; rojo: grupo sacarosa; verde: grupo stevia. Los corchetes indican las comparaciones de Tukey significativas entre tratamientos dentro de cada edad ($^{*}$~$p < 0{,}05$; $^{**}$~$p < 0{,}01$; $^{***}$~$p < 0{,}001$; $^{****}$~$p < 0{,}0001$); $n = 10$--$17$ animales por grupo.}
  \label{fig:of-rearing}
\end{figure}

La Tabla~\ref{tab:of-anova} resume los ANOVA de dos vías (edad $\times$ tratamiento); las comparaciones \textit{post hoc} de Tukey se realizaron entre los tres tratamientos dentro de cada edad.

\begin{table}[htbp]
\centering
\caption[Resumen del ANOVA de dos vías para las variables del OF]{Resumen del ANOVA de dos vías (edad $\times$ tratamiento) para las variables del OF. En negrita, los efectos significativos ($p < 0{,}05$). $\eta^2_p$: eta cuadrado parcial, proporción de la variabilidad de la variable atribuible a cada factor (tamaño del efecto).}
\label{tab:of-anova}
\footnotesize
\setlength{\tabcolsep}{2.5pt}
\begin{tabular*}{\linewidth}{@{\extracolsep{\fill}}l ccc ccc ccc@{}}
\toprule
& \multicolumn{3}{c}{\textbf{Interacción}} & \multicolumn{3}{c}{\textbf{Tratamiento}} & \multicolumn{3}{c}{\textbf{Edad}} \\
\cmidrule{2-4}\cmidrule{5-7}\cmidrule{8-10}
\textbf{Variable} & \textbf{$F$ (gl)} & \textbf{$p$} & \textbf{$\eta^2_p$} & \textbf{$F$ (gl)} & \textbf{$p$} & \textbf{$\eta^2_p$} & \textbf{$F$ (gl)} & \textbf{$p$} & \textbf{$\eta^2_p$} \\
\midrule
Distancia total & 0,06 (2, 66) & 0,94 & 0,002 & 0,90 (2, 66) & 0,41 & 0,027 & 0,31 (1, 66) & 0,58 & 0,005 \\
Zonas visitadas & 3,28 (2, 63) & \textbf{0,044} & 0,094 & 37,01 (2, 63) & \textbf{$<$0,0001} & 0,540 & 0,03 (1, 63) & 0,87 & 0,000 \\
Entradas a zona 1+2 & 1,88 (2, 57) & 0,16 & 0,062 & 91,79 (2, 57) & \textbf{$<$0,0001} & 0,763 & 3,51 (1, 57) & 0,066 & 0,058 \\
Distancia en zona 1+2 & 5,05 (2, 69) & \textbf{0,0090} & 0,128 & 38,88 (2, 69) & \textbf{$<$0,0001} & 0,530 & 3,15 (1, 69) & 0,080 & 0,044 \\
Entradas a zona 1 & 2,29 (2, 72) & 0,11 & 0,060 & 4,11 (2, 72) & \textbf{0,020} & 0,102 & 6,20 (1, 72) & \textbf{0,015} & 0,079 \\
Distancia en zona 1 & 0,94 (2, 74) & 0,39 & 0,025 & 2,80 (2, 74) & 0,067 & 0,070 & 3,85 (1, 74) & 0,053 & 0,049 \\
\textit{Rearing}, n.º de eventos & 6,67 (2, 73) & \textbf{0,0022} & 0,155 & 12,69 (2, 73) & \textbf{$<$0,0001} & 0,258 & 3,54 (1, 73) & 0,064 & 0,046 \\
\textit{Rearing}, tiempo & 2,24 (2, 72) & 0,11 & 0,059 & 8,81 (2, 72) & \textbf{0,0004} & 0,197 & 10,12 (1, 72) & \textbf{0,0022} & 0,123 \\
\bottomrule
\end{tabular*}
\end{table}

% ---------------------------------------------------------------------
\FloatBarrier
\section{Memoria de reconocimiento de objetos (NOR)}\label{sec:res-nor}

Durante la fase de muestreo (T1), el tiempo de exploración de los dos objetos idénticos reveló un efecto significativo de la edad (ANOVA de tres vías; $F_{\text{Edad}}(1, 49) = 36{,}365$; $p < 0{,}0001$; Figura~\ref{fig:nor-t1}), evidenciando que los animales juveniles presentaron un mayor tiempo de exploración en comparación con los adultos. No se observaron efectos principales significativos del tratamiento ($F_{\text{Tratamiento}}(2, 49) = 0{,}045$; $p = 0{,}9563$) ni de la posición del objeto ($F_{\text{Posición}}(1, 49) = 1{,}755$; $p = 0{,}1913$). Tampoco resultaron significativas las interacciones tratamiento $\times$ edad ($F(2, 49) = 0{,}663$; $p = 0{,}5201$), tratamiento $\times$ posición ($F(2, 49) = 0{,}260$; $p = 0{,}7719$) ni tratamiento $\times$ edad $\times$ posición ($F(2, 49) = 2{,}515$; $p = 0{,}0912$). En conjunto, estos resultados indican que la edad constituye un factor determinante del comportamiento exploratorio durante la fase de muestreo, mientras que la exposición a sacarosa o stevia no produjo diferencias significativas en el tiempo promedio de exploración de los objetos.

\begin{figure}[htbp]
  \centering
  % figura compuesta (herramientas/figuras/componer.py)
  \incluirfigura[width=\linewidth]{compuestas/Fig5_NOR_T1_compuesta.pdf}
  \caption[Exploración de los objetos en la fase de muestreo del NOR]{Tiempo de exploración de cada uno de los dos objetos idénticos durante la fase de muestreo (T1) del NOR. Barras llenas: objeto izquierdo; barras rayadas: objeto derecho. Las barras indican la media $\pm$ SEM. Azul: grupo control; rojo: grupo sacarosa; verde: grupo stevia. El corchete indica el efecto principal de la edad, que abarca a todos los grupos de cada edad ($^{****}$~$p < 0{,}0001$; ANOVA de tres vías). $n = 9$--$10$ animales por grupo.}
  \label{fig:nor-t1}
\end{figure}

Tras un intervalo de 2\,h (memoria a corto plazo, T2), se encontró una interacción edad-tratamiento significativa en la capacidad de discriminar el objeto novedoso del familiar (ANOVA de dos vías, $F_{\text{Interacción}}(2, 49) = 10{,}13$; $p = 0{,}0002$). Los grupos control juvenil y adulto mostraron una clara preferencia por explorar el objeto novedoso (tasa de exploración de 0,632 y de 0,703, respectivamente; prueba $t$ de una muestra, $t(8) = 10{,}31$ y $t(9) = 10{,}56$; $p < 0{,}0001$ en ambos casos), al igual que el grupo adulto expuesto a sacarosa. En cambio, el grupo juvenil expuesto a sacarosa exploró menos el objeto novedoso que su control de edad (prueba \textit{post hoc} de Tukey, control vs.\ sacarosa: $p = 0{,}0017$; Figura~\ref{fig:nor-tasa}A). En cuanto a la stevia, el grupo juvenil no difirió de su control (control vs.\ stevia: $p = 0{,}4613$), mientras que el grupo adulto exploró menos el objeto novedoso que su control de edad (control vs.\ stevia: $p = 0{,}0003$). Solo se detectaron diferencias por edad entre los grupos expuestos a sacarosa (juvenil vs.\ adulto: $p < 0{,}0001$); los grupos control ($p = 0{,}0591$) y stevia ($p = 0{,}2959$) no difirieron entre edades.

De manera similar, la interacción edad-tratamiento fue significativa a las 24\,h (memoria a largo plazo, T3; $F_{\text{Interacción}}(2, 49) = 11{,}67$; $p < 0{,}0001$). Los grupos control mantuvieron la preferencia por el objeto novedoso (juvenil: 0,716; $t(8) = 8{,}62$; $p < 0{,}0001$; adulto: 0,653; $t(9) = 4{,}72$; $p = 0{,}0011$). El grupo juvenil expuesto a sacarosa volvió a explorar menos el objeto novedoso que su control de edad (control vs.\ sacarosa: $p = 0{,}0008$), al igual que el grupo juvenil expuesto a stevia (control vs.\ stevia: $p = 0{,}0005$; Figura~\ref{fig:nor-tasa}B). En el grupo adulto no hubo diferencias entre tratamientos ($p > 0{,}05$). Se detectaron diferencias por edad en los grupos expuestos a sacarosa (juvenil vs.\ adulto: $p < 0{,}0001$) y a stevia ($p = 0{,}0004$), pero no en los grupos control ($p = 0{,}0965$).

Siete días después de la fase de muestreo (memoria consolidada, T4), la interacción edad-tratamiento también fue significativa ($F_{\text{Interacción}}(2, 49) = 7{,}82$; $p = 0{,}0011$). Los grupos control conservaron la preferencia por el objeto novedoso (juvenil: 0,677; $t(8) = 15{,}66$; $p < 0{,}0001$; adulto: 0,662; $t(9) = 4{,}61$; $p = 0{,}0013$). Nuevamente, el grupo juvenil expuesto a sacarosa exploró menos el objeto novedoso que su control de edad (control vs.\ sacarosa: $p = 0{,}0009$; Figura~\ref{fig:nor-tasa}C), mientras que no se registraron diferencias entre los otros grupos y sus controles de edad ($p > 0{,}05$). El grupo adulto expuesto a stevia no difirió de su control, pero exploró menos el objeto novedoso que el grupo adulto expuesto a sacarosa (sacarosa vs.\ stevia: $p = 0{,}0382$). Solo se detectaron diferencias por edad entre los grupos expuestos a sacarosa (juvenil vs.\ adulto: $p < 0{,}0001$); los grupos control ($p = 0{,}7177$) y stevia ($p = 0{,}5189$) no difirieron entre edades.

\begin{figure}[htbp]
  \centering
  % figura compuesta (herramientas/figuras/componer.py)
  \incluirfigura[width=\linewidth]{compuestas/Fig5_NOR_tasa.pdf}
  \caption[Tasa de exploración en la prueba NOR]{Tasa de exploración en la prueba NOR. (A)~T2 (2\,h). (B)~T3 (24\,h). (C)~T4 (7 días). El protocolo y los objetos de cada ensayo se muestran en la Figura~\ref{fig:nor-esquema}. Las barras indican media $\pm$ SEM; cada punto representa un animal. Azul: grupo control; rojo: grupo sacarosa; verde: grupo stevia. Se indican los resultados de las comparaciones de Tukey entre grupos ($^{*}\,p < 0{,}05$; $^{**}\,p < 0{,}01$; $^{***}\,p < 0{,}001$; $^{****}\,p < 0{,}0001$). $n = 9$--$10$ animales por grupo.}
  \label{fig:nor-tasa}
\end{figure}

En conjunto, estos resultados indican que los efectos del consumo de sacarosa y de stevia sobre la memoria de reconocimiento dependen tanto del tipo de endulzante como de la etapa del desarrollo en la que ocurre la exposición. En particular, los animales expuestos a sacarosa durante la etapa juvenil presentaron las alteraciones más pronunciadas, con déficits persistentes en la memoria a corto plazo, a largo plazo y consolidada. En contraste, los animales expuestos a sacarosa durante la adultez conservaron la capacidad de reconocer el objeto familiar y discriminarlo del novedoso en todas las fases experimentales.

Por el contrario, la exposición a stevia produjo alteraciones específicas en función de la edad y del intervalo de retención. En los animales juveniles se observó una disminución de la memoria a largo plazo, mientras que los adultos la presentaron solo en la memoria a corto plazo, sin cambios en la memoria consolidada.

Colectivamente, estos hallazgos sugieren que la exposición a sacarosa durante la etapa juvenil ejerce un efecto más amplio y persistente sobre la memoria de reconocimiento, mientras que los efectos de la stevia son más específicos y dependen de la edad de exposición y del intervalo de retención evaluado. La Tabla~\ref{tab:nor-anova} resume los ANOVA de dos vías (edad $\times$ tratamiento).

\begin{table}[htbp]
\centering
\caption[Resumen del ANOVA de dos vías para la tasa de exploración en la prueba NOR]{Resumen del ANOVA de dos vías (edad de exposición $\times$ tratamiento) para la tasa de exploración en cada demora de la prueba NOR. $\eta^2_p$: eta cuadrado parcial, proporción de la variabilidad de la variable atribuible a cada factor (tamaño del efecto).}
\label{tab:nor-anova}
\footnotesize
\setlength{\tabcolsep}{3pt}
\begin{tabularx}{\linewidth}{@{}l c c C c c C c c C@{}}
\toprule
& \multicolumn{3}{c}{\textbf{T2}} & \multicolumn{3}{c}{\textbf{T3}} & \multicolumn{3}{c}{\textbf{T4}} \\
\cmidrule(lr){2-4}\cmidrule(lr){5-7}\cmidrule(l){8-10}
\textbf{Fuente} & \textbf{$F$ (gl)} & \textbf{$p$} & \textbf{$\eta^2_p$} & \textbf{$F$ (gl)} & \textbf{$p$} & \textbf{$\eta^2_p$} & \textbf{$F$ (gl)} & \textbf{$p$} & \textbf{$\eta^2_p$} \\
\midrule
Interacción & 10,13 (2, 49) & \textbf{0,0002} & 0,293 & 11,67 (2, 49) & \textbf{$<$0,0001} & 0,323 & 7,82 (2, 49) & \textbf{0,0011} & 0,242 \\
Tratamiento & 7,89 (2, 49) & \textbf{0,0011} & 0,244 & 1,87 (2, 49) & 0,165 & 0,071 & 3,11 (2, 49) & 0,054 & 0,113 \\
Edad        & 12,74 (1, 49) & \textbf{0,0008} & 0,206 & 14,92 (1, 49) & \textbf{0,0003} & 0,233 & 9,07 (1, 49) & \textbf{0,0041} & 0,156 \\
\bottomrule
\end{tabularx}
\end{table}

% ---------------------------------------------------------------------
\FloatBarrier
\section{Separación de patrones espaciales (SLR)}\label{sec:res-slr}

La prueba SLR evalúa la capacidad de distinguir una posición novedosa de una familiar, con una demanda de separación de patrones mayor en la configuración s-SLR que en la d-SLR (Figura~\ref{fig:esquema-slr}). Los datos se informan como DI (Figura~\ref{fig:slr-di}) y como tasa de exploración equivalente (Figura~\ref{fig:slr-tasa}). Como el DI es una transformación lineal de la tasa de exploración ($\mathrm{DI} = 2\cdot\text{tasa} - 1$), ambos índices producen resultados estadísticos idénticos (mismos $F$ y $p$).

\label{sec:res-slr-exploracion}En la fase de muestreo (día 1), se analizó el tiempo de exploración de cada uno de los tres objetos idénticos (Figura~\ref{fig:slr-dia1}A) mediante un ANOVA de tres vías (edad $\times$ tratamiento $\times$ objeto). En la d-SLR, los animales exploraron los tres objetos por igual ($F_{\text{Objeto}}(2, 104) = 1{,}14$; $p = 0{,}32$), sin efectos de las interacciones ($p > 0{,}05$), del tratamiento ($F_{\text{Tratamiento}}(2, 52) = 2{,}06$; $p = 0{,}14$) ni de la edad ($F_{\text{Edad}}(1, 52) = 0{,}52$; $p = 0{,}47$). En la s-SLR, los animales exploraron más el objeto A1 que los objetos A2 y A3, próximos entre sí ($F_{\text{Objeto}}(2, 106) = 8{,}07$; $p = 0{,}0014$, con corrección de Geisser-Greenhouse), sin interacción del objeto con la edad ($p = 0{,}12$) ni con el tratamiento ($p = 0{,}46$); el tiempo de exploración dependió, además, de la interacción entre la edad y el tratamiento y del tratamiento, pero no de la edad ($F_{\text{Edad} \times \text{Tratamiento}}(2, 53) = 5{,}95$; $p = 0{,}0047$; $F_{\text{Tratamiento}}(2, 53) = 7{,}44$; $p = 0{,}0014$; $F_{\text{Edad}}(1, 53) = 0{,}16$; $p = 0{,}69$): los juveniles expuestos a sacarosa y los adultos expuestos a stevia exploraron más los objetos que sus controles de edad (prueba \textit{post hoc} de Tukey, control vs.\ sacarosa en juveniles: $p < 0{,}01$; control vs.\ stevia en adultos: $p < 0{,}01$). Además, los animales expuestos a sacarosa y a stevia recorrieron mayor distancia que los controles durante la fase de muestreo en la s-SLR (ANOVA de dos vías sobre la distancia recorrida por minuto; d-SLR: $F_{\text{Interacción}}(2, 48) = 0{,}02$; $p = 0{,}98$; $F_{\text{Tratamiento}}(2, 48) = 5{,}62$; $p = 0{,}006$; $F_{\text{Edad}}(1, 48) = 3{,}31$; $p = 0{,}075$; prueba \textit{post hoc} de Tukey, control vs.\ tratado: $p > 0{,}05$; s-SLR: $F_{\text{Interacción}}(2, 51) = 0{,}16$; $p = 0{,}85$; $F_{\text{Tratamiento}}(2, 51) = 14{,}57$; $p < 0{,}0001$; $F_{\text{Edad}}(1, 51) = 7{,}37$; $p = 0{,}009$; control vs.\ sacarosa: juvenil, $p < 0{,}05$; adulto, $p < 0{,}01$; control vs.\ stevia: juvenil, $p > 0{,}05$; adulto, $p < 0{,}01$; Figura~\ref{fig:slr-dia1}B).

\begin{figure}[htbp]
  \centering
  % figura compuesta (herramientas/figuras/componer.py)
  \incluirfigura[width=\linewidth]{compuestas/Fig5_SLR_dia1_compuesta.pdf}
  \caption[Exploración de los objetos y distancia recorrida en la fase de muestreo de la SLR]{Fase de muestreo (día 1) de la SLR, en las configuraciones d-SLR y s-SLR. (A)~Tiempo de exploración de cada uno de los tres objetos idénticos. A1: objeto cuya posición se mantiene en la fase de prueba (posición familiar); A2 y A3: objetos entre cuyas posiciones se ubica, en la fase de prueba, el objeto en posición novedosa. (B)~Distancia recorrida por minuto de sesión. Las barras indican la media $\pm$ SEM. Azul: grupo control; rojo: grupo sacarosa; verde: grupo stevia. Los corchetes indican las comparaciones de Tukey significativas con el control dentro de cada edad ($^{*}$~$p < 0{,}05$; $^{**}$~$p < 0{,}01$). $n = 9$--$10$ animales por grupo y configuración (en B, $n = 6$--$10$: no hay registro de distancia de una tanda de controles juveniles ni de un control adulto).}
  \label{fig:slr-dia1}
\end{figure}

% La figura de la exploración total día 1 vs. día 2 (fig:slr-exploracion, Fig5_exploracion_SLR.pdf)
% y su frase se quitaron a pedido de Sol (23 oct. 2026, comentario C89); quedan en la etiqueta tesina-20261022.

En la configuración d-SLR, la discriminación de la posición novedosa dependió del tratamiento (ANOVA de dos vías, $F_{\text{Interacción}}(2, 52) = 2{,}04$; $p = 0{,}140$; $F_{\text{Tratamiento}}(2, 52) = 4{,}58$; $p = 0{,}015$; $F_{\text{Edad}}(1, 52) = 0{,}22$; $p = 0{,}642$). Los grupos control juvenil y adulto exploraron preferentemente el objeto ubicado en la posición novedosa (prueba $t$ de una muestra, $p = 0{,}001$ y $p = 0{,}003$, respectivamente). El grupo juvenil expuesto a sacarosa mostró un DI menor que su control de edad (prueba \textit{post hoc} de Tukey, control vs.\ sacarosa: $p = 0{,}009$), mientras que en el grupo adulto no hubo diferencias entre tratamientos ($p > 0{,}50$); la comparación entre edades dentro del grupo sacarosa no alcanzó la significancia ($p = 0{,}067$). Los grupos expuestos a stevia no difirieron de sus controles (juvenil: $p = 0{,}870$; adulto: $p > 0{,}50$), y el grupo juvenil expuesto a stevia mostró un DI mayor que el grupo juvenil expuesto a sacarosa (sacarosa vs.\ stevia: $p = 0{,}027$; Figura~\ref{fig:slr-di}A).

En la configuración s-SLR, de mayor demanda de separación de patrones, la discriminación también dependió del tratamiento (ANOVA de dos vías, $F_{\text{Interacción}}(2, 53) = 0{,}53$; $p = 0{,}591$; $F_{\text{Tratamiento}}(2, 53) = 6{,}67$; $p = 0{,}003$; $F_{\text{Edad}}(1, 53) = 0{,}08$; $p = 0{,}779$). Los grupos control volvieron a discriminar la posición novedosa por encima del nivel de azar (juvenil: $p < 0{,}001$; adulto: $p = 0{,}007$). En esta configuración, los animales expuestos a sacarosa mostraron un DI menor que sus controles en ambas edades (prueba \textit{post hoc} de Tukey, control vs.\ sacarosa: juvenil, $p = 0{,}041$; adulto, $p = 0{,}032$; Figura~\ref{fig:slr-di}B), sin diferencias entre edades dentro del tratamiento. Dado que los residuos de este modelo se apartaron de la normalidad (Shapiro-Wilk, $p = 0{,}006$), estas diferencias se contrastaron también con la prueba no paramétrica de Mann-Whitney, que las confirmó en ambas edades (control vs.\ sacarosa: juvenil, $p = 0{,}002$; adulto, $p = 0{,}031$); la homogeneidad de varianzas se mantuvo (Brown-Forsythe, $p = 0{,}16$), al igual que en la d-SLR (Shapiro-Wilk, $p = 0{,}48$; Brown-Forsythe, $p = 0{,}30$). Los grupos expuestos a stevia no difirieron de sus controles (control vs.\ stevia: juvenil, $p = 0{,}971$; adulto, $p = 0{,}293$), y la diferencia entre el grupo juvenil expuesto a stevia y el expuesto a sacarosa no alcanzó la significancia ($p = 0{,}069$).

\begin{figure}[htbp]
  \centering
  % figura compuesta (herramientas/figuras/componer.py)
  \incluirfigura[width=\linewidth]{compuestas/Fig5_SLR_DI.pdf}
  \caption[DI en la prueba SLR]{DI en la prueba SLR. (A)~Configuración d-SLR. (B)~Configuración s-SLR. Las barras indican media $\pm$ SEM. Azul: grupo control; rojo: grupo sacarosa; verde: grupo stevia. Se indican los resultados de las comparaciones de Tukey entre grupos ($^{*}\,p < 0{,}05$; $^{**}\,p < 0{,}01$). $n = 9$--$10$ animales por grupo.}
  \label{fig:slr-di}
\end{figure}

\begin{figure}[htbp]
  \centering
  % figura compuesta (herramientas/figuras/componer.py)
  \incluirfigura[width=\linewidth]{compuestas/Fig5_SLR_tasa.pdf}
  \caption[Tasa de exploración en la prueba SLR]{Tasa de exploración en la prueba SLR. (A)~Configuración d-SLR. (B)~Configuración s-SLR. Las barras indican media $\pm$ SEM. Azul: grupo control; rojo: grupo sacarosa; verde: grupo stevia. Se indican los resultados de las comparaciones de Tukey entre grupos ($^{*}\,p < 0{,}05$; $^{**}\,p < 0{,}01$), estadísticamente idénticos a los del DI (Figura~\ref{fig:slr-di}) por tratarse de una transformación lineal. $n = 9$--$10$ animales por grupo.}
  \label{fig:slr-tasa}
\end{figure}

En conjunto, estos resultados indican: (1) la exposición a sacarosa durante la etapa juvenil redujo la discriminación de la posición novedosa en ambas configuraciones, mientras que la exposición durante la adultez la redujo solo en la s-SLR, la de mayor demanda de separación de patrones; y (2) los animales expuestos a stevia no difirieron de los controles en ninguna condición. La Tabla~\ref{tab:slr-anova} resume los ANOVA de dos vías (edad $\times$ tratamiento) de cada configuración, y la Figura~\ref{fig:slr-tukey} resume las diferencias entre grupos junto con sus intervalos de confianza.

\begin{figure}[htbp]
  \centering
  \incluirfigura[width=\linewidth]{Fig5_estimacion_Tukey_SLR.png}
  \caption[Diferencias de medias del DI entre grupos (prueba de Tukey)]{Diferencias de medias del DI entre grupos (prueba de Tukey) con su intervalo de confianza del 95\%, para las configuraciones d-SLR (A) y s-SLR (B). Cada fila corresponde a una comparación; una diferencia positiva indica un DI mayor en el primer grupo de la comparación. En rojo se indican las comparaciones significativas ($p < 0{,}05$). La línea punteada indica ausencia de diferencia.}
  \label{fig:slr-tukey}
\end{figure}

\begin{table}[htbp]
\centering
\caption[Resumen del ANOVA de dos vías para el DI en la prueba SLR]{Resumen del ANOVA de dos vías (edad de exposición $\times$ tratamiento) para el DI en cada configuración de la prueba SLR. $\eta^2_p$: eta cuadrado parcial, proporción de la variabilidad de la variable atribuible a cada factor (tamaño del efecto).}
\label{tab:slr-anova}
\footnotesize
\begin{tabularx}{\linewidth}{@{}l C c c C c c@{}}
\toprule
& \multicolumn{3}{c}{\textbf{d-SLR}} & \multicolumn{3}{c}{\textbf{s-SLR}} \\
\cmidrule(lr){2-4}\cmidrule(l){5-7}
\textbf{Fuente de variación} & \textbf{$F$ (gl)} & \textbf{$p$} & \textbf{$\eta^2_p$} & \textbf{$F$ (gl)} & \textbf{$p$} & \textbf{$\eta^2_p$} \\
\midrule
Interacción & 2,04 (2, 52) & 0,140 & 0,073 & 0,53 (2, 53) & 0,591 & 0,020 \\
Tratamiento & 4,58 (2, 52) & \textbf{0,015} & 0,150 & 6,67 (2, 53) & \textbf{0,003} & 0,201 \\
Edad        & 0,22 (1, 52) & 0,642 & 0,004 & 0,08 (1, 53) & 0,779 & 0,001 \\
\bottomrule
\end{tabularx}
\end{table}

% ---------------------------------------------------------------------
\FloatBarrier
\section{Neurogénesis hipocampal (PCNA y DCX)}\label{sec:res-neurogenesis}

La proliferación celular y la neurogénesis en el GD se evaluaron mediante la inmunodetección de PCNA, marcador de células en proliferación, y de DCX, marcador de neuroblastos y neuronas inmaduras, en los animales de los grupos control y sacarosa (Figuras~\ref{fig:ihc} y~\ref{fig:ihc-fotos}).

Los animales expuestos a sacarosa presentaron menos células PCNA$^{+}$ que los controles en ambas edades (ANOVA de dos vías sin término de interacción, $F_{\text{Tratamiento}}(1, 77) = 17{,}98$; $p < 0{,}0001$; $F_{\text{Edad}}(1, 77) = 3{,}29$; $p = 0{,}074$; prueba \textit{post hoc} de Tukey, control vs.\ sacarosa: juvenil, $p = 0{,}0004$; adulto, $p = 0{,}0004$; Figura~\ref{fig:ihc}A). Dentro de cada tratamiento, los grupos juvenil y adulto no difirieron entre sí ($p = 0{,}27$).

En cambio, el número de células DCX$^{+}$ no difirió entre grupos ($F_{\text{Interacción}}(1, 76) = 3{,}22$; $p = 0{,}077$; $F_{\text{Tratamiento}}(1, 76) = 0{,}96$; $p = 0{,}33$; $F_{\text{Edad}}(1, 76) = 0{,}03$; $p = 0{,}87$; Figura~\ref{fig:ihc}B), y ninguna comparación de Tukey resultó significativa (control vs.\ sacarosa: juvenil, $p = 0{,}94$; adulto, $p = 0{,}21$).

Las células doblemente marcadas (PCNA$^{+}$/DCX$^{+}$), que corresponden a neuroblastos en proliferación, disminuyeron en los animales expuestos a sacarosa en ambas edades (ANOVA de dos vías, $F_{\text{Interacción}}(1, 76) = 0{,}13$; $p = 0{,}72$; $F_{\text{Tratamiento}}(1, 76) = 38{,}35$; $p < 0{,}0001$; $F_{\text{Edad}}(1, 76) = 1{,}19$; $p = 0{,}28$; prueba \textit{post hoc} de Tukey, control vs.\ sacarosa: juvenil, $p < 0{,}0001$; adulto, $p = 0{,}0005$; Figura~\ref{fig:ihc}C). En este marcador, la dispersión de los valores fue marcadamente menor en los grupos expuestos a sacarosa (Brown-Forsythe, $p < 0{,}0001$).

\begin{figure}[htbp]
  \centering
  % figura compuesta (herramientas/figuras/componer.py)
  \incluirfigura[width=\linewidth]{compuestas/Fig5_IHC_compuesta.pdf}
  \caption[IHC del GD (PCNA y DCX)]{IHC del GD en ratas expuestas a agua (control) o a sacarosa durante la ventana juvenil o adulta. (A)~Células PCNA$^{+}$. (B)~Células DCX$^{+}$. (C)~Células PCNA$^{+}$/DCX$^{+}$. Las barras indican la media $\pm$ SEM. Azul: grupo control; rojo: grupo sacarosa. Se indican las comparaciones de Tukey significativas entre control y sacarosa dentro de cada edad ($^{***}\,p < 0{,}001$; $^{****}\,p < 0{,}0001$; $N = 4$, $n = 5$).}
  \label{fig:ihc}
\end{figure}

\begin{figure}[htbp]
  \centering
  % herramientas/figuras/ihc_fotos_horizontal.py (28 oct. 2026: 2 filas x 4 grupos)
  \incluirfigura[width=\linewidth]{Fig5_IHC_fotos_horizontal.png}
  \caption[Fotomicrografías representativas del GD (PCNA y DCX)]{Fotomicrografías representativas de la inmunorreactividad para PCNA (verde) y DCX (rojo) en el GD de animales control y expuestos a sacarosa durante la ventana juvenil o adulta. Las flechas blancas señalan células PCNA$^{+}$/DCX$^{+}$, y los círculos muestran una ampliación de la zona indicada. Objetivo de 20$\times$; barra de escala: \SI{100}{\micro\metre}.}
  \label{fig:ihc-fotos}
\end{figure}

En conjunto, estos resultados indican: (1) la exposición a sacarosa redujo la proliferación celular en el GD en ambas edades; (2) redujo el número de neuroblastos en proliferación (PCNA$^{+}$/DCX$^{+}$) en ambas edades; y (3) no modificó el número total de células DCX$^{+}$. Como las células PCNA$^{+}$/DCX$^{+}$ son las que todavía se dividen dentro del linaje neuronal, su reducción sugiere que la sacarosa disminuye la generación de nuevas neuronas en el GD, aun cuando el número total de neuronas inmaduras no cambie.

% ---------------------------------------------------------------------
\FloatBarrier
\section{Análisis multivariado de la prueba SLR}\label{sec:res-ml}

Para evaluar si las métricas de la prueba SLR definen perfiles conductuales propios de cada grupo, se analizaron en conjunto cinco variables por animal: el tiempo de exploración de cada objeto en la fase de muestreo (A1, A2 y A3), la exploración total en la fase de prueba y el DI. Las dos primeras componentes principales (CP) explicaron el 69,3\% de la varianza (Figura~\ref{fig:ml-pca}). La primera (CP1, 49,7\%) reunió los cuatro tiempos de exploración, con cargas similares (0,46 a 0,51), y representa el nivel general de exploración de los objetos; la segunda (CP2, 19,6\%) estuvo dominada por el DI (carga 0,91) y representa la discriminación de la posición novedosa. Ambas dimensiones fueron, por lo tanto, casi independientes: explorar más no implicó discriminar mejor.

Los animales juveniles expuestos a sacarosa y los adultos expuestos a stevia presentaron los valores más altos de CP1, es decir, exploraron más los objetos en ambas sesiones, mientras que los controles juveniles presentaron los más bajos (ANOVA de dos vías, $F_{\text{Interacción}}(2, 111) = 10{,}34$; $p = 0{,}0001$; $F_{\text{Tratamiento}}(2, 111) = 8{,}20$; $p = 0{,}0005$; $F_{\text{Edad}}(1, 111) = 0{,}06$; $p = 0{,}81$; Figura~\ref{fig:ml-pca}A). En cambio, los animales expuestos a sacarosa presentaron los valores más bajos de CP2 en ambas edades ($F_{\text{Interacción}}(2, 111) = 0{,}62$; $p = 0{,}54$; $F_{\text{Tratamiento}}(2, 111) = 6{,}38$; $p = 0{,}0024$; $F_{\text{Edad}}(1, 111) = 0{,}38$; $p = 0{,}54$), lo que reproduce en el espacio multivariado el déficit de discriminación de la sacarosa. Los grupos expuestos a stevia se ubicaron en CP2 junto a los controles. Ninguna de las dos componentes difirió entre la cohorte de 2026 y las de 2024--2025 (prueba $t$; CP1: $p = 0{,}37$; CP2: $p = 0{,}32$).

El clustering de Ward no recuperó los grupos experimentales. La partición con mayor coeficiente de silueta (dos clusters; silueta = 0,37) separó un subconjunto de 23 animales con exploración elevada en ambas sesiones, integrado por animales de cuatro de los seis grupos, del resto de la muestra, y la concordancia entre los clusters y los grupos fue prácticamente nula en todas las particiones evaluadas (ARI $\leq 0{,}10$ entre 2 y 6 clusters). Esto indica que los animales no se organizan espontáneamente en perfiles conductuales separados según la exposición.

El clasificador Random Forest predijo el grupo experimental por encima del azar (exactitud balanceada = $0{,}30 \pm 0{,}08$, frente a un azar de 0,17; prueba de permutación, $p = 0{,}006$; Figura~\ref{fig:ml-rf}A). Los grupos mejor clasificados fueron los juveniles expuestos a sacarosa (54\% de aciertos) y los controles juveniles (50\%), mientras que los juveniles expuestos a stevia fueron los peor clasificados (6\%) y se distribuyeron entre los demás grupos. El DI fue la variable de mayor importancia (Figura~\ref{fig:ml-rf}B). Al restringir la clasificación a dos tratamientos, el modelo distinguió a los controles tanto de los animales expuestos a sacarosa (exactitud balanceada = $0{,}67 \pm 0{,}09$; $p = 0{,}010$) como de los expuestos a stevia ($0{,}67 \pm 0{,}10$; $p = 0{,}002$), pero no discriminó los tres tratamientos a la vez ($0{,}46 \pm 0{,}08$, frente a un azar de 0,33; $p = 0{,}080$). Al incluir la distancia recorrida en la fase de muestreo, la exactitud para los seis grupos aumentó a $0{,}39 \pm 0{,}10$ ($p = 0{,}002$), pero esa misma variable permitió identificar la cohorte de 2026 con una exactitud de $0{,}83 \pm 0{,}08$ ($p = 0{,}002$), frente a $0{,}62 \pm 0{,}09$ sin ella ($p = 0{,}066$), lo que justificó excluirla del análisis principal.

\begin{figure}[htbp]
  \centering
  \incluirfigura[width=\linewidth]{Fig5_ML_pca.pdf}
  \caption[PCA de la prueba SLR]{PCA de cinco métricas de la prueba SLR por animal: tiempo de exploración de los objetos A1, A2 y A3 en la fase de muestreo (día 1), exploración total en la fase de prueba (día 2) y DI. (A)~Puntuaciones de cada animal en las dos primeras componentes (CP1 y CP2; círculos llenos: juveniles; triángulos vacíos: adultos) y media $\pm$ SEM de cada grupo (símbolos grandes). Cada punto es un animal; los animales con perfiles parecidos quedan cerca. Valores altos de CP1 indican más exploración de los objetos; valores altos de CP2, mayor DI. (B)~Cargas de cada variable en CP1 y CP2: cuánto contribuye cada variable a cada componente (las barras largas indican las variables que la definen; el signo, si aumenta o disminuye con ella). Azul: grupo control; rojo: grupo sacarosa; verde: grupo stevia. $n = 117$ animales (18--20 por grupo, d-SLR y s-SLR combinadas).}
  \label{fig:ml-pca}
\end{figure}

\begin{figure}[htbp]
  \centering
  \incluirfigura[width=\linewidth]{Fig5_ML_rf.pdf}
  \caption[Clasificación de los grupos experimentales con Random Forest]{Clasificación de los seis grupos experimentales con Random Forest a partir de las cinco métricas de la prueba SLR. (A)~Matriz de confusión: proporción de animales de cada grupo real asignados a cada grupo predicho en los conjuntos de prueba (validación cruzada estratificada de 5 particiones, repetida 10 veces); la diagonal corresponde a los aciertos, y un clasificador que respondiera al azar acertaría en 0,17 de los casos (1 de 6). Juv.: juvenil; Ad.: adulto; C: control; S: sacarosa; St: stevia. (B)~Importancia de cada variable, estimada como la disminución de la exactitud balanceada al mezclar al azar sus valores en los conjuntos de prueba (media $\pm$ DE de las 50 particiones): cuanto mayor la caída, más depende el clasificador de esa variable. $n = 117$ animales.}
  \label{fig:ml-rf}
\end{figure}

En conjunto, estos resultados indican: (1) las métricas de la SLR se organizan en dos dimensiones casi independientes, el nivel de exploración y la discriminación de la posición novedosa; (2) la sacarosa se distingue del control por una menor discriminación en ambas edades, mientras que la stevia se distingue del control por su nivel de exploración, pero no por su discriminación; y (3) ningún grupo constituye un perfil conductual separado del resto, y los juveniles expuestos a stevia son los que menos se diferencian de los demás grupos.

% =====================================================================
\chapter{Discusión}\label{cap:discusion}
% Reescrita el 19 oct. 2026 con el estilo de las discusiones de la Dra. Kruse
% (Kruse et al. 2019; Coirini et al. 2022). 20 oct. 2026: correcciones de Sol
% (docx "metodos resultados discusion 20261019 (1)"): la sacarosa es el control
% positivo y el foco es la stevia; la stevia tambien produce alteraciones (perfil
% ansiolitico y deficits de memoria de reconocimiento), no se la presenta como
% inocua; SLR y neurogenesis en una misma seccion; sin seccion de Limitaciones;
% conclusion de Sol. Las afirmaciones sobre otros trabajos se verificaron contra
% PubMed (Tsan et al. evaluaron la memoria durante el consumo y el efecto se vio
% con los tres edulcorantes juntos, no con la stevia sola).
% =====================================================================

\section{Síntesis de los hallazgos}\label{sec:disc-sintesis}

En este trabajo se evaluaron los efectos a largo plazo del consumo de un extracto de stevia durante la etapa juvenil o durante la adultez y se los comparó con los de la sacarosa al 10\,\%, cuyos efectos persistentes sobre la conducta y la memoria ya fueron descriptos por el laboratorio \cite{kruse2019sucrose,coirini2022sucrose}. Los animales expuestos a stevia exploraron más las zonas centrales del \textit{open field} en ambas edades, lo que indica un perfil de tipo ansiolítico, y presentaron alteraciones de la memoria de reconocimiento que dependieron de la edad de exposición ---a largo plazo en los juveniles y a corto plazo en los adultos---, sin afectar la memoria consolidada ni la separación de patrones espaciales. En cambio, los animales expuestos a sacarosa durante la etapa juvenil exploraron menos las zonas centrales del \textit{open field}, no lograron discriminar el objeto novedoso en ninguna de las demoras del NOR y discriminaron peor la posición novedosa en ambas configuraciones de la SLR, y los expuestos a sacarosa en ambas edades presentaron menos neuroblastos en proliferación en el giro dentado. Estos resultados respaldan la hipótesis de que la stevia no reproduce los efectos adversos de la sacarosa (\S\ref{sec:hipotesis}), pero muestran que no está exenta de efectos a largo plazo: produce alteraciones de la conducta distintas y más sutiles que las de la sacarosa.

\section{Consumo de líquidos}\label{sec:disc-consumo}

Durante el tratamiento, los animales juveniles consumieron, en relación con su peso corporal, más del doble de solución endulzada que los adultos, tanto de sacarosa como de stevia (\S\ref{sec:res-consumo-tratamiento}). Sin embargo, el índice de preferencia no presentó diferencias significativas entre edades, lo que indica que esta mayor ingesta relativa no estuvo acompañada de una mayor preferencia por el sabor dulce. En consecuencia, a pesar de presentar una preferencia comparable por las soluciones endulzadas, los animales juveniles consumieron cantidades considerablemente superiores de sacarosa y de stevia por unidad de peso corporal que los adultos. Asimismo, la ingesta de stevia superó la ingesta diaria admisible de glucósidos de esteviol, de 4\,mg/kg de peso corporal por día expresada como equivalentes de esteviol \cite{efsa2010steviol}, y este exceso fue más pronunciado durante la etapa juvenil. Según la composición del extracto (Tabla~\ref{tab:extracto}, con los factores de conversión a esteviol de 0,40 para el esteviósido y de 0,33 para los rebaudiósidos A y C \cite{efsa2010steviol}), la bebida al 4\,\% contenía alrededor de 0,08\,mg/mL de equivalentes de esteviol; con un consumo medio de 0,24 y 0,09\,mL/g de peso corporal por día, los animales juveniles ingirieron alrededor de 20\,mg/kg/día (unas cinco veces la ingesta diaria admisible) y los adultos, alrededor de 7\,mg/kg/día (1,8 veces).

Los animales prefirieron la sacarosa más que la stevia en ambas edades. Este resultado es compatible con la propuesta de que el sabor dulce disociado del aporte calórico se procesa de forma distinta en los circuitos de recompensa \cite{yunker2020nns}, y con el hallazgo de que el consumo de stevia aumenta la expresión de $\Delta$FosB, un marcador de activación sostenida, en el núcleo accumbens, el caudado-putamen y la corteza prefrontal, regiones del sistema de recompensa; en ese trabajo, la cría en un ambiente enriquecido redujo tanto el consumo de stevia como ese aumento de $\Delta$FosB \cite{salinasvelarde2021fosb}. Una menor preferencia por la solución no calórica puede considerarse un aspecto favorable de la stevia, ya que sugiere un menor potencial para promover un consumo excesivo; de hecho, los animales consumieron menos stevia que sacarosa en ambas ventanas de edad (\S\ref{sec:res-consumo-tratamiento}).

Una vez retirados los edulcorantes, solo los animales expuestos durante la etapa juvenil modificaron su consumo de agua (\S\ref{sec:res-consumo-lavado}). Los juveniles expuestos a sacarosa bebieron menos agua que su control de edad durante las primeras dos semanas, aunque la diferencia no alcanzó la significancia, y luego se aproximaron a los valores del control. Este comportamiento se asocia con frustración, como fue descripto previamente por el laboratorio \cite{kruse2019sucrose}. Por el contrario, los juveniles expuestos a stevia bebieron menos agua que su control de edad durante los primeros nueve días, y más entre los días 11 y 25. Por otro lado, los adultos expuestos a stevia o a sacarosa no presentaron diferencias respecto de su control de edad en ningún momento, lo que refuerza la mayor sensibilidad de la ventana juvenil.

\section{Conducta tipo ansiosa y exploración en el \textit{open field}}\label{sec:disc-of}

En el \textit{open field}, los animales expuestos a sacarosa durante la etapa juvenil ingresaron menos veces a las zonas centrales, recorrieron en ellas menor distancia, visitaron menos zonas y realizaron menos \textit{rearing}, mientras que los expuestos a sacarosa durante la adultez no difirieron de su control (\S\ref{sec:res-of}), lo que indica que el consumo temprano de una bebida azucarada promueve un perfil ansioso a largo plazo, como fue descripto anteriormente \cite{kruse2019sucrose}.

A pesar de que no se registraron diferencias en la distancia total recorrida entre los grupos, los animales que bebieron stevia visitaron más zonas que sus controles de edad y que los grupos expuestos a sacarosa. Principalmente, los animales expuestos a stevia ingresaron más veces a las zonas centrales y recorrieron en ellas mayor distancia que el resto de los grupos, independientemente de la edad en la que se produjo el consumo.

Estudios previos del laboratorio mostraron que el consumo temprano de stevia (decocción al 0,5\,\% v/v) no produce cambios en la conducta tipo ansiosa de ratas macho juveniles cuando se evalúa a corto plazo, al finalizar la exposición \cite{coirini2023stevia,marchena2026stevia} \confirmar{coautores y revista del manuscrito de Marchena et~al.}. Esto sugiere que la abstinencia de la bebida endulzada con stevia, algunas alteraciones neurobiológicas, o ambas, serían responsables de los cambios a largo plazo en la conducta de tipo ansioso. Lo mismo ocurre con la sacarosa, cuyos efectos sobre la conducta tipo ansiosa evaluados a corto plazo \cite{sanchezhuerta2022sucrose,rey2024metformin} difieren de los descriptos a largo plazo \cite{kruse2019sucrose}. De manera similar, en ratas macho con acceso a stevia durante la adolescencia y evaluadas al finalizar el consumo, la stevia aumentó la locomoción en el \textit{open field} \cite{rico2024stevia}, mientras que en este trabajo, a largo plazo, no modificó la distancia total recorrida y aumentó la exploración de las zonas centrales. Por otra parte, resultados previos del laboratorio demostraron que la exposición temprana al mismo extracto acuoso en ratas hembra juveniles disminuye las conductas de tipo ansioso en el laberinto en cruz elevado \cite{kruse2025stevia}. Estos hallazgos sugieren la existencia de diferencias sexuales en la respuesta conductual al extracto. En este contexto, sería de interés investigar si los efectos ansiolíticos observados en las hembras persisten a largo plazo y si la edad de exposición influye en su magnitud y persistencia.

Estos resultados están en línea con los que demuestran efectos ansiolíticos de la stevia en ratas diabéticas \cite{khakpai2023stevia}. Los animales expuestos a stevia realizaron, además, menos \textit{rearing} (los juveniles respecto del control y los adultos en el tiempo dedicado), una conducta exploratoria vertical asociada a la novedad y al hipocampo \cite{lever2006rearing}; sin embargo, este cambio no se acompañó de una menor exploración de la arena, sino de una mayor exploración del centro. En conjunto, estos resultados indican que la stevia no reproduce el perfil ansioso de la sacarosa a largo plazo y que, por el contrario, se asocia con una mayor exploración de los espacios abiertos.

\section{Memoria de reconocimiento de objetos}\label{sec:disc-nor}

En el NOR, los animales expuestos a sacarosa durante la etapa juvenil presentaron una menor exploración del objeto novedoso que sus controles de edad en todos los intervalos de retención evaluados: a las 2\,h (memoria a corto plazo), a las 24\,h (memoria a largo plazo) y a los 7 días (memoria consolidada). En contraste, los animales expuestos a sacarosa durante la adultez conservaron la capacidad de discriminar entre el objeto familiar y el novedoso en los tres intervalos de retención (\S\ref{sec:res-nor}), como fue descripto previamente \cite{coirini2022sucrose}.

Por el contrario, la exposición a stevia no alteró la memoria consolidada en ninguna de las dos edades. Sin embargo, los animales del grupo stevia presentaron dificultades para discriminar los objetos según la edad de exposición: los adultos presentaron alteraciones en la memoria a corto plazo, y los juveniles, en la memoria a largo plazo.

De manera similar a lo observado en el \textit{open field}, estudios previos del laboratorio indican que el consumo agudo de stevia durante la etapa juvenil no produce alteraciones en la memoria de localización de objetos (prueba NOL), una memoria dependiente del hipocampo \cite{coirini2023stevia,marchena2026stevia}. Además, estos animales presentaron modificaciones en la expresión de calretinina (CR), un marcador de determinadas poblaciones de interneuronas GABAérgicas, en la corteza prefrontal medial \cite{marchena2026stevia}. Las interneuronas que expresan parvalbúmina (PV) y CR son importantes reguladoras de la actividad de los circuitos neuronales y desempeñan funciones relevantes en procesos cognitivos como la memoria y la sincronización de las redes neuronales \cite{bosch2015amygdala,ognjanovski2017pv}. En este sentido, estudios previos demostraron que la disminución de la expresión de PV en el hipocampo y en la corteza prefrontal tras el consumo de sacarosa se asocia con alteraciones de la memoria y de las conductas de tipo ansioso \cite{reichelt2015adolescent,kruse2019sucrose,xu2018sucrose,reichelt2019hfhs}. Considerando estos antecedentes, sería de interés investigar si las modificaciones en la expresión de CR tras el consumo juvenil de stevia persisten hasta la edad adulta y si se relacionan con las alteraciones de la memoria de reconocimiento observadas a largo plazo.

Son escasos los trabajos que evaluaron los efectos de la stevia sobre la memoria a largo plazo, es decir, después de finalizado el consumo. La mayoría la evaluó mientras los animales aún consumían el edulcorante, con resultados dispares. Morales-Ríos et~al.\ encontraron que el consumo de stevia durante 18 semanas alteró la memoria a largo plazo en el NOR de ratas macho adultas \cite{moralesrios2022sweeteners}. En cambio, Tsan et~al., que administraron edulcorantes de bajas calorías (acesulfamo de potasio, sacarina o stevia) desde la etapa juvenil y evaluaron la memoria dependiente del hipocampo durante el período de consumo, describieron un deterioro al analizar los tres edulcorantes en conjunto, pero no al analizar la stevia por separado \cite{tsan2022lcs}. A largo plazo, solo se describió que el consumo materno de stevia durante la gestación y la lactancia se asocia con un peor desempeño de memoria en la descendencia macho adulta \cite{delagarza2022maternal}. Este trabajo aporta, por lo tanto, evidencia de que el consumo de stevia en la etapa juvenil o en la adultez produce alteraciones de la memoria de reconocimiento que persisten después de finalizado el consumo, aunque más limitadas que las de la sacarosa.

A diferencia de esos estudios, que administraron la stevia en dosis cercanas a la ingesta diaria admisible, aquí se administró un extracto acuoso de hojas que contiene, además de glucósidos de esteviol, compuestos fenólicos y flavonoides con actividad antioxidante (Tabla~\ref{tab:extracto}). La stevia redujo los marcadores de estrés oxidativo cerebral en ratas alimentadas con una dieta alta en grasas y sacarosa \cite{pozdnyakova2025stevia}, y los glucósidos de esteviol presentan actividad antioxidante \cite{prata2017glycosides}, revirtieron el deterioro de la memoria inducido por escopolamina \cite{singh2010stevioside} y modularon la plasticidad sináptica hipocampal en un modelo de desorden metabólico \cite{chavushyan2017stevia}. Estas propiedades podrían contribuir a que la stevia no produjera un deterioro sostenido de la memoria de reconocimiento.

\section{Separación de patrones espaciales y neurogénesis}\label{sec:disc-slr}\label{sec:disc-neurogenesis}

En la SLR, los animales expuestos a sacarosa durante la etapa juvenil discriminaron peor la posición novedosa en ambas configuraciones, mientras que los expuestos durante la adultez lo hicieron solo en la s-SLR, la configuración de mayor demanda de separación de patrones (\S\ref{sec:res-slr}). Este patrón concuerda con el de Reichelt et~al., quienes encontraron que las ratas con acceso diario a sacarosa discriminaron la nueva ubicación de un objeto cuando la separación espacial era grande, pero no cuando era pequeña \cite{reichelt2016sucrose}. También concuerda con la alteración selectiva de la memoria de lugar descripta tras una exposición breve a azúcar \cite{beilharz2014place,beilharz2016liquid} y con el deterioro de la separación de patrones asociado a una menor neurogénesis tras una alteración nutricional temprana \cite{perezgarcia2016malnutrition}. Que en los adultos el déficit apareciera solo en la configuración más exigente sugiere que la exposición adulta a sacarosa compromete específicamente la separación de patrones, mientras que la exposición juvenil produce una alteración más amplia de la memoria espacial.

Los animales expuestos a stevia, en cambio, discriminaron la posición novedosa de la misma forma que los controles en ambas configuraciones y en ambas edades, y en la d-SLR los juveniles expuestos a stevia discriminaron mejor que los expuestos a sacarosa. Este resultado indica que la exposición a stevia no altera la separación de patrones espaciales. La exploración de los objetos durante la fase de muestreo fue mayor en los grupos expuestos a los edulcorantes; sin embargo, como los animales expuestos a sacarosa exploraron más, y no menos, que los controles, el déficit de discriminación de la sacarosa no se explica por una exploración insuficiente durante la codificación.

En concordancia con los resultados de la SLR, los animales expuestos a sacarosa presentaron alteraciones en el proceso de neurogénesis. La exposición a sacarosa redujo el número de células PCNA$^{+}$ en el giro dentado, con una diferencia significativa en el grupo adulto, y redujo el número de neuroblastos en proliferación (PCNA$^{+}$/DCX$^{+}$) independientemente de la edad en la que se produjo el consumo, sin modificar el número total de células DCX$^{+}$ (\S\ref{sec:res-neurogenesis}). Estos resultados indican que la sacarosa afecta la etapa proliferativa del linaje neuronal \cite{kempermann2015neurogenesis} y que esa alteración persiste 25 días después de finalizado el consumo. Reichelt et~al.\ también encontraron menos células PCNA$^{+}$ en el giro dentado de ratas macho jóvenes con acceso diario a sacarosa al 10\,\% durante 28 días, aunque las evaluaron mientras aún consumían la solución y encontraron, además, menos células DCX$^{+}$ \cite{reichelt2016sucrose}. En este trabajo, en cambio, el número total de células DCX$^{+}$ no se modificó, como en el trabajo de Xu y Reichelt, en el que la sacarosa sola no modificó la DCX del giro dentado \cite{xu2018sucrose}; esto podría indicar que, a largo plazo, la alteración se restringe a la etapa proliferativa. Además, en ratas hembra, el consumo de sacarosa al 25\,\% durante doce semanas a partir del destete redujo la expresión hipocampal de \textit{Neurod1}, un gen vinculado a la diferenciación neuronal \cite{maniam2016sugar}. En cambio, Sánchez-Huerta et~al.\ encontraron un aumento de la proliferación en la zona subgranular del giro dentado ventral tras el consumo de sacarosa iniciado en la adolescencia tardía \cite{sanchezhuerta2022sucrose}, lo que sugiere que el efecto sobre la proliferación depende de la edad de inicio de la exposición.

La reducción de los neuroblastos en proliferación se observó en ambas edades, al igual que el déficit de la sacarosa en la s-SLR. La separación de patrones depende de la neurogénesis del giro dentado: cuando la neurogénesis disminuye, los animales dejan de discriminar ubicaciones cercanas, que exigen mayor separación de patrones, pero siguen discriminando ubicaciones alejadas \cite{reichelt2016sucrose,reichelt2021slr}. Por eso, Reichelt et~al.\ propusieron que la menor proliferación inducida por la sacarosa subyace al déficit en la s-SLR \cite{reichelt2016sucrose}. El déficit en la s-SLR y la reducción de los neuroblastos en proliferación constituyen, entonces, dos maneras de observar un mismo efecto de la sacarosa: la alteración de la neurogénesis hipocampal. En cambio, no está claro hasta qué punto la memoria de reconocimiento de objetos depende de la neurogénesis; de hecho, los adultos expuestos a sacarosa presentaron menos proliferación sin alteraciones en el NOR, lo que sugiere que esta prueba depende en menor medida de este proceso. La cuantificación de estos marcadores en los animales expuestos a stevia permitirá determinar si la preservación de la separación de patrones observada en esos grupos se acompaña de una neurogénesis conservada.

\section{Análisis multivariado}\label{sec:disc-ml}

El análisis multivariado de la SLR (\S\ref{sec:res-ml}) mostró que las métricas de la prueba se organizan en dos dimensiones casi independientes: el nivel de exploración de los objetos y la discriminación de la posición novedosa. En la dimensión de la discriminación, los grupos expuestos a sacarosa se ubicaron por debajo de los controles en ambas edades, mientras que los expuestos a stevia se ubicaron junto a los controles, lo que confirma con un método independiente los resultados de los ANOVA. Que el DI haya sido la variable más importante para el clasificador refuerza que el déficit de la sacarosa es un déficit de memoria y no de exploración. La stevia se distinguió del control por su nivel de exploración, que fue mayor en los adultos expuestos a stevia, lo que concuerda con la mayor exploración de las zonas centrales del \textit{open field} (\S\ref{sec:disc-of}). Ni el clustering ni el clasificador encontraron grupos conductuales separados, y la exactitud del clasificador, superior al azar pero modesta, es la esperable para una intervención nutricional: en los trabajos que aplicaron Random Forest al fenotipado conductual, los grupos diferían por manipulaciones farmacológicas o genéticas de gran efecto \cite{lukovikov2026cylinder,geros2020tracking}.

\section{Ventana etaria}\label{sec:disc-edad}

Los resultados de este trabajo muestran que la etapa juvenil es más vulnerable a los efectos de la sacarosa que la adultez. En el NOR, la interacción edad $\times$ tratamiento fue significativa en las tres memorias evaluadas (memoria a corto plazo, a largo plazo y consolidada; Tabla~\ref{tab:nor-anova}), y la memoria se alteró principalmente con la exposición juvenil a sacarosa, aunque la exposición adulta a stevia también alteró la memoria a corto plazo; en el \textit{open field}, solo la exposición juvenil a sacarosa aumentó la conducta tipo ansiosa; y en la SLR, la exposición juvenil a sacarosa afectó ambas configuraciones, mientras que la adulta afectó solo la más exigente.

En conjunto, estos resultados sugieren que la etapa juvenil constituye una ventana de especial vulnerabilidad frente a los efectos del consumo de endulzantes. Este período se caracteriza por una intensa reorganización y maduración de los circuitos neuronales, que incluye procesos de plasticidad sináptica, refinamiento de las conexiones neuronales y maduración de los sistemas de neurotransmisión, fundamentales para el desarrollo de las funciones cognitivas y emocionales. En este contexto, la exposición a determinados estímulos durante esta etapa podría interferir con dichos procesos y desencadenar modificaciones neurobiológicas persistentes, que podrían mantenerse hasta la edad adulta. Si bien las alteraciones asociadas al consumo de stevia fueron más limitadas que las observadas tras la exposición a sacarosa, se requieren estudios adicionales para esclarecer los mecanismos neurobiológicos involucrados. Asimismo, resulta necesario profundizar en la evaluación de la seguridad del consumo de stevia durante etapas tempranas del desarrollo, particularmente en relación con sus posibles efectos sobre la maduración cerebral y las funciones cognitivas.

\section{En resumen}\label{sec:disc-integracion}

Los resultados de este trabajo muestran que el consumo ilimitado de una bebida endulzada con sacarosa durante la etapa juvenil afecta al cerebro en desarrollo de forma distinta y más duradera que si el consumo ocurre en la adultez. En particular, la exposición juvenil a sacarosa aumentó la conducta tipo ansiosa, alteró la memoria de reconocimiento a corto plazo, a largo plazo y consolidada, deterioró la separación de patrones espaciales y redujo la proliferación de neuroblastos en el giro dentado. La exposición a un extracto de stevia, en cambio, no reprodujo estos efectos: los animales conservaron la memoria consolidada y la separación de patrones, aun cuando la demanda cognitiva era mayor. Sin embargo, los animales expuestos a stevia presentaron un perfil conductual de tipo ansiolítico, acompañado de alteraciones en la memoria de reconocimiento a corto y a largo plazo, independientemente de la edad en la que se produjo el consumo.

En conclusión, los resultados obtenidos, junto con la evidencia aportada por estudios previos, sugieren que el consumo crónico de bebidas azucaradas durante etapas tempranas de la vida puede afectar negativamente el desarrollo y la maduración del sistema nervioso central, con consecuencias que podrían persistir hasta la edad adulta. Por su parte, el consumo de stevia se asoció con un perfil de efectos diferente, caracterizado por conductas de tipo ansiolítico y alteraciones más moderadas de la memoria de reconocimiento, observadas tanto en animales juveniles como adultos. Si bien estos efectos fueron de menor magnitud que los asociados al consumo de sacarosa, su presencia en ambas etapas de la vida, sumada a los efectos sobre la fertilidad descriptos en ratas hembra \cite{kruse2025stevia}, pone de manifiesto que la stevia no está exenta de efectos. En este sentido, resulta fundamental continuar investigando sus posibles consecuencias neurobiológicas y cognitivas a largo plazo, así como profundizar en la evaluación de la seguridad de su consumo en las distintas etapas de la vida.

% =====================================================================
\chapter{Perspectivas futuras}\label{cap:perspectivas}
% Plantilla ITBA: opcional, máximo 1 carilla. Puede incluir objetivos de
% máxima no alcanzados. Borrador del 21 oct. 2026: incluye, como trabajo a
% futuro, las limitaciones del diseño que la Dra. Kruse sacó de la Discusión
% (a confirmar con ella).
% =====================================================================

Los resultados de este trabajo abren distintas líneas para continuarlo:

\begin{itemize}
  \item \textbf{Neurogénesis en los grupos expuestos a stevia.} Cuantificar PCNA y DCX en los animales expuestos a stevia permitirá determinar si la preservación de la separación de patrones espaciales se acompaña de una neurogénesis conservada. Conviene, además, registrar el animal de origen de cada corte, de modo que el análisis pueda realizarse por animal y no por corte.
  \item \textbf{Mecanismos.} Evaluar si las modificaciones en la expresión de calretinina y parvalbúmina descriptas tras el consumo agudo de stevia persisten a largo plazo y se relacionan con las alteraciones de la memoria de reconocimiento, y estudiar el papel de la abstinencia de la bebida endulzada en los cambios de la conducta tipo ansiosa.
  \item \textbf{Hembras.} Extender el estudio a ratas hembra, en las que el mismo extracto redujo la conducta tipo ansiosa y alteró la fertilidad, para establecer si esos efectos persisten a largo plazo y si dependen de la edad de exposición.
  \item \textbf{Dosis.} Dado que la ingesta superó los valores diarios de referencia, evaluar concentraciones menores del extracto, cercanas a la ingesta diaria admisible de glucósidos de esteviol.
  \item \textbf{Diseño de las cohortes.} Incluir un grupo control en cada cohorte experimental y ampliar el número de animales expuestos a stevia en el \textit{open field}.
  \item \textbf{Análisis multivariado.} Registrar las métricas de \anymaze{} identificadas por animal también en el \textit{open field} y en el NOR, para extender el PCA, el clustering y el Random Forest a toda la batería conductual e integrar en un mismo análisis las distintas pruebas, junto con el peso corporal.
\end{itemize}

% =====================================================================
\chapter{Conclusiones}\label{cap:conclusiones}
% Plantilla ITBA: máximo 2 carillas; indicar si se cumplieron los
% objetivos y si la hipótesis fue comprobada o refutada.
% Borrador del 21 oct. 2026 (pendiente de revisión de la Dra. Kruse): una
% conclusión por asunto, estado de cada objetivo y de la hipótesis. No agrega
% resultados nuevos: resume los de la sección 5 con el enfoque de la Discusión.
% =====================================================================

En este trabajo se evaluaron los efectos a largo plazo del consumo de una bebida endulzada con un extracto de stevia durante la etapa juvenil o la adultez, y se los comparó con los de la sacarosa, utilizada como control positivo por sus efectos ya descriptos. De los resultados obtenidos se desprenden las siguientes conclusiones.

\textbf{Consumo.} Los animales juveniles consumieron, en relación con su peso corporal, más del doble de solución endulzada que los adultos, sin una mayor preferencia por el sabor dulce, y la ingesta de ambos \resaltar{endulzantes} superó los valores diarios de referencia.

\textbf{Conducta tipo ansiosa.} El consumo de stevia, en cualquiera de las dos edades, se asoció a largo plazo con una mayor exploración de las zonas centrales del \textit{open field}, sin cambios en la actividad locomotora, es decir, con un perfil de tipo ansiolítico. La sacarosa, en cambio, aumentó la conducta tipo ansiosa solo cuando se consumió durante la etapa juvenil.

\textbf{Memoria de reconocimiento.} La stevia alteró la memoria de reconocimiento de objetos de manera dependiente de la edad: a largo plazo en los animales expuestos durante la etapa juvenil y a corto plazo en los expuestos durante la adultez, sin afectar la memoria consolidada. La sacarosa consumida durante la etapa juvenil alteró los tres tipos de memoria evaluados.

\textbf{Separación de patrones espaciales y neurogénesis.} La stevia no alteró la separación de patrones espaciales en ninguna de las dos edades. La sacarosa la deterioró en ambas edades en la configuración de mayor demanda, y también en la de menor demanda cuando se consumió durante la etapa juvenil, y redujo en ambas edades la proliferación celular y el número de neuroblastos en proliferación en el \resaltarb{GD}: dos expresiones de un mismo efecto sobre la neurogénesis hipocampal.

\textbf{Análisis multivariado.} Las métricas de la prueba SLR se organizaron en dos dimensiones casi independientes, la exploración de los objetos y la discriminación de la posición novedosa. Los animales expuestos a stevia se ubicaron junto a los controles en la dimensión de la discriminación, y el clasificador distinguió a cada \resaltar{endulzante} del control por encima del azar, aunque ningún grupo constituyó un perfil conductual separado del resto.

\section*{Cumplimiento de los objetivos}

Se cumplieron los objetivos de mínima referidos al registro del consumo, a la conducta tipo ansiosa, a la memoria de reconocimiento y a la separación de patrones espaciales. \resaltar{La} cuantificación de PCNA y DCX se completó en los grupos control y sacarosa, pero no en los grupos expuestos a stevia. Los objetivos de máxima se cumplieron sobre la prueba SLR: el PCA y el clustering de Ward mostraron que los animales no se agrupan espontáneamente según su exposición, y el Random Forest predijo el grupo de cada animal por encima del azar, con el índice de discriminación como la variable de mayor importancia.

\section*{Hipótesis}

La hipótesis se cumplió solo en parte. La stevia no reprodujo los efectos adversos de la sacarosa sobre la memoria consolidada y la separación de patrones espaciales, pero no resultó inocua: produjo, en ambas edades, un perfil de tipo ansiolítico y alteraciones de la memoria de reconocimiento, de menor magnitud que las de la sacarosa. Su efecto sobre la neurogénesis hipocampal queda por determinar.

En conclusión, el consumo de stevia en etapas tempranas de la vida o en la adultez no está exento de efectos a largo plazo sobre el comportamiento, aunque estos son más sutiles que los de la sacarosa. Estos resultados indican que la seguridad de la stevia como alternativa al azúcar, especialmente durante el desarrollo, requiere una evaluación más profunda.


\chapter*{Referencias bibliográficas}
\printbibliography[heading=none]

\end{document}
